% Options for packages loaded elsewhere
\PassOptionsToPackage{unicode}{hyperref}
\PassOptionsToPackage{hyphens}{url}
\documentclass[11pt]{extarticle}
\usepackage{kotex}
\usepackage{xcolor}
\usepackage[margin=18mm]{geometry}
\usepackage{amsmath,amssymb}
\setcounter{secnumdepth}{-\maxdimen} % remove section numbering
\usepackage{iftex}
\ifPDFTeX
  \usepackage[T1]{fontenc}
  \usepackage[utf8]{inputenc}
  \usepackage{textcomp} % provide euro and other symbols
\else % if luatex or xetex
  \usepackage{unicode-math} % this also loads fontspec
  \defaultfontfeatures{Scale=MatchLowercase}
  \defaultfontfeatures[\rmfamily]{Ligatures=TeX,Scale=1}
\fi
\usepackage{lmodern}
\ifPDFTeX\else
  % xetex/luatex font selection
\fi
% Use upquote if available, for straight quotes in verbatim environments
\IfFileExists{upquote.sty}{\usepackage{upquote}}{}
\IfFileExists{microtype.sty}{% use microtype if available
  \usepackage[]{microtype}
  \UseMicrotypeSet[protrusion]{basicmath} % disable protrusion for tt fonts
}{}
\makeatletter
\@ifundefined{KOMAClassName}{% if non-KOMA class
  \IfFileExists{parskip.sty}{%
    \usepackage{parskip}
  }{% else
    \setlength{\parindent}{0pt}
    \setlength{\parskip}{6pt plus 2pt minus 1pt}}
}{% if KOMA class
  \KOMAoptions{parskip=half}}
\makeatother
\usepackage{longtable,booktabs,array,ragged2e}
\usepackage{enumitem}
\newcounter{none} % for unnumbered tables
\usepackage{calc} % for calculating minipage widths
% Correct order of tables after \paragraph or \subparagraph
\usepackage{etoolbox}
\makeatletter
\patchcmd\longtable{\par}{\if@noskipsec\mbox{}\fi\par}{}{}
\makeatother
% Allow footnotes in longtable head/foot
\IfFileExists{footnotehyper.sty}{\usepackage{footnotehyper}}{\usepackage{footnote}}
\makesavenoteenv{longtable}
\setlength{\emergencystretch}{3em} % prevent overfull lines
\providecommand{\tightlist}{%
  \setlength{\itemsep}{0pt}\setlength{\parskip}{0pt}}
\usepackage{bookmark}
\IfFileExists{xurl.sty}{\usepackage{xurl}}{} % add URL line breaks if available
\urlstyle{same}
\hypersetup{
  pdftitle={zkDPP Decision Log},
  pdfauthor={Protocol Design Rationale},
  hidelinks,
  pdfcreator={LaTeX via pandoc}}

\newcommand{\List}[1]{\mathrm{List}_{#1}}
\newcommand{\Map}[1]{\mathrm{Map}_{#1}}
\newcolumntype{L}[1]{>{\RaggedRight\arraybackslash}p{#1}}
\DeclareMathOperator{\HashToField}{HashToField}
\DeclareMathOperator{\PoseidonTwo}{Poseidon2}

\title{zkDPP Decision Log}
\author{Protocol Design Rationale}
\date{2026-07-16}

\begin{document}
\maketitle
\tableofcontents
\newpage

\section{DPP Protocol Decision Log -
요약}\label{dpp-protocol-decision-log---uxc694uxc57d}

이 표는 현재 유효한 결론을 빠르게 확인하기 위한 Summary다. 각 결론의 검토 과정,
반려안과 변경 이력은 뒤의 상세 기록에서 확인한다.

\begin{center}\rule{0.5\linewidth}{0.5pt}\end{center}

\subsection{빠른 결론}\label{uxbe60uxb978-uxacb0uxb860}

{\def\LTcaptype{none} % do not increment counter
\small
\renewcommand{\arraystretch}{1.3}
\setlength{\tabcolsep}{5pt}
\begin{longtable}[]{@{}L{0.24\textwidth}L{0.56\textwidth}L{0.12\textwidth}@{}}
\toprule\noalign{}
\textbf{고민거리} & \textbf{결론} & \textbf{상태} \\
\midrule\noalign{}
\endhead
\bottomrule\noalign{}
\endlastfoot
기존 GPT 대화는 어디에 있는가? &
\texttt{/Users/bighim/Downloads/ChatGPT-DPP\ 수식\ 정리.md}에 있다. &
확인됨 \\
이번 decision이 Recall만인가? & 아니다. proof stack, Sahai positioning,
DPP data model, note commitment, event relation, provenance/privacy,
Recall이 모두 포함된다. & 확정 \\
구현 stack은 무엇인가? & Go, gnark v0.15.0, PLONK, BLS12-381, KZG,
Poseidon2, Solidity와 Foundry를 사용한다. & 확정 \\
문서에서 Field는 무엇을 의미하는가? & 별도 설명이 없으면 BLS12-381 scalar field
$\mathbb{F}_r$를 뜻한다. Base field $\mathbb{F}_q$는 curve point 좌표와 pairing
구현에서 사용한다. & 확정 \\
Poseidon2 parameter는 무엇을 사용하는가? & gnark v0.15.0이 의존하는
gnark-crypto v0.20.1의 BLS12-381 default instance를 그대로 사용한다.
Width 2, full rounds 6, partial rounds 50, S-box $x^5$와
Merkle--Damg\aa rd construction이다. & POC 확정 \\
Groth16/aggregation/recursion은 main인가? & 아니다. optimization 또는
future work다. & 확정 \\
Sahai와의 관계는 무엇인가? & Sahai-style document graph를 DPP-specific
private state payload와 event-specific relation으로 확장한다. & 확정 \\
Product information을 어떻게 나누는가? & \texttt{DocumentInfo}는
$\HashToField(\mathit{DocumentInfoBytes})$로 binding하고,
$q$와 $\mathbf{s}$는 note payload의 private witness로 둔다. & 확정 \\
Hash-only 방식인가? & 아니다. hash-only integrity model이 아니라
Attribute Linking / Consistency Check model이다. & 확정 \\
\texttt{DocumentHash}의 역할은 무엇인가? & 반복 검증하지 않는 일반
product information을 binding한다. & 확정 \\
POC \texttt{DocumentInfo}에는 무엇을 넣는가? &
\(\mathit{DocumentInfo}=(\mathit{ProductName},\mathit{LotId},\mathit{Unit})\)으로 시작한다.
더 필요한 product identity/version/type 정보는 \texttt{cm}이 아니라
\texttt{DocumentInfo}에 넣는다. & POC 확정 \\
\texttt{DocumentHash}는 어떻게 만드는가? &
각 문자열을 NFC 정규화하고 UTF-8로 encode한다. UTF-8 byte 길이를 unsigned
32-bit big-endian으로 먼저 기록해 연결한 뒤
gnark-crypto BLS12-381 \texttt{fr.Hash}로 scalar field 원소 하나를 만든다.
Layout version과 hash 용도 구분값은 넣지 않고, API의 DST에는 빈 byte string을
전달한다. & POC 확정 \\
\texttt{C\_state}를 유지하는가? & 유지하지 않는다. 비례분배를 검증하려면
어차피 circuit witness와 commitment를 link해야 하므로 Pedersen
commit-and-prove overhead가 크다. & 확정 \\
Pedersen을 사용하는가? & POC 설계에서는 제거한다. 외부 homomorphic check
대신 Split/Merge/Process relation을 circuit 내부에서 검증한다. & 확정 \\
note commitment는 무엇을 쓰는가? & POC에서는 다음 flat layout을 사용한다.
\[
cm=\PoseidonTwo(\mathit{DocumentHash},q,s_1,\ldots,s_k,addr,o)
\]
& POC 확정 \\
State에는 어떤 속성을 넣는가? & \texttt{quantity}는 별도 field로
분리하고, \texttt{StateVector}에는 \texttt{total\_kg},
\texttt{carbon\_kgCO2e}, \texttt{recycled\_kg}처럼 quantity에
비례분배되는 amount-like attributes를 둔다. & 확정 \\
\texttt{cm}에 product type, version, schema id를 넣는가? & 넣지 않는다.
Product identity나 versioning에 필요한 값은 \texttt{DocumentInfo}에
넣고, \texttt{cm}에는 consistency-critical note payload만 flat하게
넣는다. & 확정 \\
percentage, grade, certificate 같은 값은 State인가? & 직접 State에
넣기보다 amount 값으로 변환하거나 \texttt{DocumentInfo}, policy, 별도
predicate/category로 둔다. & 확정 \\
소수점 단위는 어떻게 다루는가? & \texttt{quantity}는 integer count로
두고, \texttt{total\_kg}, \texttt{carbon\_kgCO2e},
\texttt{recycled\_kg}는 nano scale $10^9$ fixed-point
integer로 encode한다. & POC 확정 \\
Entry는 어떤 값 범위를 검사하는가? & $0<q<2^{64}$와
$0\leq s_k<2^{64}$를 검사한다. Entry quantity는 0을 허용하지 않지만
$s_k=0$은 허용하며, 별도의 initial product semantic invariant를 추가하지 않는다. & POC 확정 \\
Transfer/Split에서 fixed-point 값이 나누어떨어지지 않으면? &
deterministic first-output remainder assignment를 사용한다. non-first output을
내림 계산하고 첫 번째 output이 residual을 받아 coordinate별 합을 정확히 보존한다.
Transfer의 첫 output은 $rv$, Split의 첫 output은 $cm_{out,1}$이다. & POC 확정 \\
Split output이 0인지 검사하는가? & 검사하지 않는다. $q_1>0$ 또는 $q_2>0$ 같은
nonzero policy는 POC 핵심 relation에 필요하지 않다. unsigned 64-bit range와
$q_{in}=q_1+q_2$만 적용한다. & POC 확정 \\
Split의 전체 합은 어떻게 확인하는가? & Pedersen 외부 homomorphic check는
사용하지 않고 SplitCircuit 내부의 first-output residual relation으로 정확히
보존한다. & 확정 \\
기본 DPP event는 무엇인가? & Entry, Transfer, Merge, Split, Process,
Exit이다. & 확정 \\
Transfer는 어떤 구조인가? & recall-default Transfer다.
$cm_{in}\rightarrow(rv,cm_C)$ 구조다.
receiver \texttt{cm}은 만들지 않고, sender change note는 만든다. &
확정 \\
Recall/Proceed는 DPP event인가? & 독립 supply-chain event가 아니라
Transfer pending state의 resolution event다. & 확정 \\
\texttt{cm} 안에 RecallVoucher를 넣는가? & 넣지 않는다. & 확정 \\
\texttt{cm\_pending}은 쓰는가? & 쓰지 않는다. & 확정 \\
\texttt{rv}는 어떻게 만드는가? & 다음과 같이 만든다.
\[
\begin{aligned}
rv=H(&\mathit{DocumentHash},q_T,\mathbf{s}_T,\\
     &addr_{Sender},addr_{Receiver},\Delta_{epoch},o_{rv})
\end{aligned}
\]
& 확정 \\
Zcash식 blind transfer처럼 같은 \texttt{cm}을 유지하는가? & 아니다.
Transfer의 sender change portion은 별도 \texttt{cm\_change}로 만든다.
Provenance event 공개 여부와는 별도의 note-state 결정이다. & 확정 \\
\texttt{rvnf}는 어떻게 만드는가? & $rvnf=H(rv,o_{rv})$로 만든다.
sender와 receiver만 아는 $o_{rv}$가 필요하다. & 확정 \\
pending transfer history는 어떻게 관리하는가? &
$EPOCH\_SIZE=600$초로 두고
$currentEpoch=\lfloor block.timestamp/EPOCH\_SIZE\rfloor$를 계산한 뒤
$\Map{rv}[rv]=currentEpoch+\Delta_{epoch}$만 둔다. 완성된 deadline을
transaction input으로 받지 않는다. status
enum은 두지 않는다. & 확정 \\
resolve 중복 방지는 어떻게 하는가? & $\List{rvnf}$의 존재 검사로 막는다.
Solidity에서는 bool mapping으로 표현할 수 있다. & 확정 \\
Recall/Proceed timing은 어떻게 되는가? & Recall은
$currentEpoch<deadline$인 동안만 가능하다. $\Delta_{epoch}=0$이면 Recall은
처음부터 불가능하다. Proceed는 Recall/Proceed로
resolve되지 않았다면 deadline 이후에도 가능하다. & 확정 \\
누가 $\Delta_{epoch}$을 선택하는가? & Transfer를 생성하는 Sender가 선택한다.
최소값은 0이고 protocol policy상 maximum은 두지 않는다. Voucher hash에 binding되는
public Field Element이지만 Circuit에서 별도 bit range check를 하지 않는다. Generated
verifier가 canonical Field 표현을 검사하고, 검증된 값으로 계산한
$currentEpoch+\Delta_{epoch}$은 \texttt{uint256} 범위 안에 있다. & 확정 \\
POC Contract 함수는 누가 호출하는가? & 공급망 Event와 Recall resolution Event는 누구나
제출할 수 있고, note/voucher 권한은 Circuit relation으로 검증한다.
\texttt{RegisterUser(addr,pkOwn)}는 호출자의 \texttt{msg.sender}를 편의용 metadata에
연결할 뿐 uniqueness나 $addr=H(pkOwn)$을 Contract에서 검사하지 않는다. 같은 EOA의
재등록은 최신 값으로 덮어쓰고 registration event는 emit하지 않는다. & POC 간소화 \\
logical List와 Map은 Solidity에서 어떻게 표현하는가? & Field Element는
\texttt{uint256}, membership List는 \texttt{mapping(uint256 => bool)},
$\Map{rv}$는 \texttt{mapping(uint256 => uint256)}으로 구현한다. 전체 set을 순회하는
dynamic array는 두지 않는다. EOA 등록은
\texttt{mapping(address => UserRegistration)}으로 구현하며 value는 $addr$와
$pkOwn$만 저장한다. & POC 확정 \\
proof public input의 순서는 무엇인가? & Protocol Specification의 $\mathbf{x}$ tuple에
적힌 왼쪽에서 오른쪽 순서가 canonical 순서다. 각 Field Element는 Solidity
\texttt{uint256} word로 전달한다. & POC 확정 \\
어떤 event를 현재 구현하는가? & Provenance graph edge event는 보류한다.
Merkle witness 획득을 위해 $MT$와 $rvMT$ append의 leaf, zero-based index와
new root는 event로 공개한다. & 부분 확정 \\
privacy claim은 무엇인가? & 현재 확정 범위는 DocumentInfo/State/witness
hiding이다. Product graph privacy는 provenance schema 결정 뒤 확정한다. & 부분 확정 \\
Process relation은 어떻게 잡는가? & 각 input/output StateVector를
$\mathbf{s}_{in,i},\mathbf{s}_{out,j}$로 쓰고, 그 aggregate를
$\mathbf{s}_{in},\mathbf{s}_{out}$으로 구분한다. 별도 State 행렬 기호는 두지 않는다.
Semantic delta는 $\mathbf{P}=(-p_{totalLoss},+p_{carbonAdded},+p_{recycledDelta})$의 State prefix를
사용한다. $\ell=3$ profile의 public input은
$P_{pub}=(p_{totalLoss},p_{carbonAdded},p_{recycledDelta})$이며
$p_{recycledDelta}$ 적용식과 recycled aggregate 보존식을 함께 검증한다. 이 두 식으로
$p_{recycledDelta}=0$이 도출되며 별도 0 assertion은 두지 않는다. Circuit은
$\mathbf{s}_{out}=\mathbf{s}_{in}+\mathbf{P}$에 대응하는 coordinate별 relation을 검증한다.
Active output은 $1\leq q_{out,j}<2^{64}$이며 quantity가 0인 output은 $n$에서 제외한다. &
POC 확정 \\
Process에서 recipe는 무엇인가? & POC에서는 recipe를
\texttt{StateAllocationMatrix\ A} 또는 \texttt{A}를 결정하는 rule로
단순화한다. on-chain recipe registry는 두지 않는다. 논문에서는 \texttt{A}를
process recipe의 state-allocation component로 표현할 수 있다. & POC 확정 \\
Policy/Recipe 기반 Process를 이번 POC에 포함하는가? & 포함하지 않는다.
Recipe가 없는 generic Process는 공개된 $P,A$에 대한 State accounting consistency만
증명한다. 사후 감사를 줄이려면 Material과 Quantity transition까지 강제하는
Recipe/Policy가 필요하지만, Recipe별 circuit은 공정을 노출하고 hidden Recipe
registry는 승인 주체와 data availability의 신뢰 문제를 추가한다. 이 방향은 추후
연구로 기록하고 현재 POC는 기존 public $P,A$ Process를 유지한다. & POC 제외/추후 고려 \\
Process의 output별 consistency는 어떻게 검증하는가? & fixed-point
allocation denominator는 constant $D=10^9$으로 둔다. $A$는 public input이다.
output 2부터 $n$까지 floor/remainder relation을
검증하고 output 1이 모든 State residual을 받도록 aggregate를 정확히 보존한다.
$\sum_j A_{j,k}=D$도 circuit 내부에서 검증한다. Client와 transaction은 active
$n\times\ell$ row-major $A$만 전달하고, Client와 Contract가 각각 fixed
$a\times\ell$ verifier statement로 zero-padding한다. & POC 확정 \\
Process에서 quantity를 강제하는가? & Process에서는 quantity
transition을 강제하지 않는다. input quantity는 기존 note에서 검증된 range를 계승하고,
output quantity는 unsigned 64-bit range와 commitment binding을 검사한다.
추후 필요하면 별도 QuantityAllocationVector
relation을 추가한다. & POC 확정 \\
Process에서 Pedersen homomorphic check를 쓰는가? & 쓰지 않는다.
Process는 allocation ratio, aggregate constraint, output commitment
binding을 fixed-arity ProcessCircuit 내부에서 직접 검증한다. & 설계
원칙 \\
Process POC의 최대 arity는 무엇인가? & 기능 확인용 기본값은
$M_{\max}=N_{\max}=3$이다. Benchmark는 $M_{\max}=N_{\max}=a$,
$a\in\{1,2,3\}$인 별도 circuit을 State 3, Depth 32에서 만들고
1-to-1, 2-to-2, 3-to-3만 측정한다.
비대칭 arity는 제외한다. & POC 확정 \\
Process active slot은 어떻게 표현하는가? & Public input으로 active count $m,n$만
전달한다. 별도 public/private enable vector는 두지 않는다. Contract는 active 배열 길이에서
$m,n$을 계산하고, fixed circuit은 $m,n$에 binding된 algebraic gating expression을 사용한다.
각 input membership path는 자신의 $rt_i$에 binding한다. & POC 확정 \\
첫 구현 milestone은 어디까지인가? & State 3, Depth 32에서 Entry와 정확한
3-to-3 Process만 end-to-end로 구현한다. $m,n$은 public input에 유지하지만 둘 다
3인지 circuit에서 검사한다. 최대 3-slot selector와 나머지 Event는 후속 구현이다.
& POC 확정 \\
State 길이와 Process benchmark 조합은? & StateVector 길이
$\ell\in\{0,1,2,3\}$을 Depth 32에서 prefix layout으로 측정한다. Process arity
1-to-1, 2-to-2, 3-to-3은 $\ell=3$에서만 수행한다. 이후 실험은 Poseidon2만 사용하고
모든 축의 Cartesian product는 만들지 않는다. & POC 확정 \\
SHA-to-Field 실험을 계속하는가? & 아니다. EVM Tree gas와 circuit overhead를 확인한
일회성 비교 결과와 코드는 보존하지만, 이후 구현과 benchmark는 Poseidon2만 사용한다.
& 확정 \\
POC address model은 어떻게 둘 것인가? &
$usk=sk_{own}$, $pk_{own}=H(sk_{own})$, $addr=H(pk_{own})$,
$upk=(addr,pk_{own})$로 둔다. & POC 확정 \\
최종 Address/identity privacy model은 확정되었는가? & 아니다. POC에서는
단순 address를 쓰고, one-time address/identity privacy는 이후 설계로
둔다. & 미확정 \\
Merge와 Split arity는 무엇인가? & Merge는 2-to-1, Split은 1-to-2로
고정하고 반복 composition한다. & POC 확정 \\
Finalized note nullifier는 어떻게 만드는가? & 별도 문자열 tag 없이
$nf=H(sk_{own},cm)$로 만든다. & POC 확정 \\
BLS12-381 Solidity verifier가 실행되는가? & gnark v0.15.0 export와 Foundry
v1.7.1 Prague EVM의 valid/invalid proof test를 통과했다. & 실험 확인 \\
DocumentHash를 circuit에서 재계산하는가? & 아니다. circuit은 Event 사이의
DocumentHash relation만 검사하고, receiver가 DocumentInfo로부터 off-chain에서
재계산한다. & 확정 \\
Note에 owner address를 저장하는가? & 저장하지 않는다. Note는
\texttt{(cm, DocumentHash, q, StateVector, o)}이고 holder는 등록된 address를
별도로 알고 사용한다. & 확정 \\
새 commitment에는 누구의 address를 넣는가? & 해당 output의 새 holder address를
넣는다. change는 sender, Proceed는 receiver, Recall은 직전 sender의 address를 쓴다.
& 확정 \\
Quantity와 State의 range는 무엇인가? & POC에서 모두 unsigned 64-bit다.
Quantity는 integer이고 State는 nano scale fixed-point다. 기존 Tree 객체의 input range는
closed-state invariant로 계승하고, 새 숫자 상태를 계산하는 output에서 검사한다. & POC 확정 \\
Transfer change quantity가 0이면 어떻게 하는가? & $q_C=0$이어도
\texttt{cm\_change}를 항상 생성하며 change State도 0이어야 한다. & POC 확정 \\
Finalized note와 Recall Voucher를 어디에 저장하는가? & $cm$은 $MT$, $rv$는
$rvMT$라는 서로 다른 append-only Merkle tree에 저장한다. & 확정 \\
Merkle tree 세부 parameter도 확정되었는가? & 기본 POC는 Depth 32로 고정하고
Depth scaling은 수행하지 않는다. Empty leaf 0, 재귀 default node,
모든 과거 root 허용과 Azeroth식 node storage로 확정했다. Client는 append event의
zero-based index를 Tree별 Contract view 함수에 전달해 current root와 path를 조회한다.
& POC 확정 \\
몇 개의 proving/verifying key를 만드는가? & Entry부터 Exit까지 8개 Relation별로
각각 PLONK proving/verifying key를 만든다. & POC 확정 \\
POC KZG SRS는 어떻게 만드는가? & \texttt{unsafekzg.NewSRS}를 사용한다.
각 Relation의 CCS에서 canonical/Lagrange SRS pair를 생성한다. Production SRS 출처는
별도 결정으로 남긴다. & POC 확정 \\
Verifier는 어떻게 배포하는가? & 8개 Relation별 generated \texttt{PlonkVerifier}를
독립 배포하고, main zkDPP contract가 constructor에서 받은 8개 address를 Relation ID에
binding한다. 공통 ABI는 \texttt{Verify(bytes,uint256[])}다. & POC 확정 \\
Setup artifact는 어떻게 관리하는가? &
\texttt{circuits/profiles/state-0}부터 \texttt{state-3}까지 네 profile folder를 두되
공통 Relation source는 복사하지 않는다. Generated artifact와 Solidity는
\texttt{state-<l>/depth-<d>/<relation-variant>/} 아래에 두고 manifest checksum으로
같은 setup 결과에 binding한다. & POC 확정 \\
Foundry는 어떻게 실행하는가? & 공식 Docker image
\texttt{ghcr.io/foundry-rs/foundry:v1.7.1}과 Prague EVM을 사용한다. & POC 확정 \\
새 note마다 key와 address를 새로 만드는가? & 아니다. POC에서는 등록된
$pk_{own}$과 $addr$을 재사용하고, 각 신규 commitment의 opening $o$만 새로 뽑는다.
& POC 확정 \\
Note와 voucher secret은 어떻게 전달하는가? & 필요한 $o$, $o_{rv}$와 opening은
off-chain으로 전달한다. Public-key encryption scheme은 현재 POC와 논문 범위에서
정의하지 않는다. & 범위 확정 \\
Entry에서 address를 새로 계산하는가? & 아니다. Client는 이미 생성하고 등록한
address를 사용해 최초 $cm$을 만든다. & 확정 \\
Entry의 기능은 무엇인가? & 공급망에 아직 존재하지 않는 Product를 최초 finalized
note로 등록하는 $\varnothing\rightarrow cm_{new}$ Event다. & 확정 \\
canonical 문서는 무엇인가? & Protocol Specification, Development
Specification과 이 Decision Log의 세 TeX 문서다. & 확정 \\
\end{longtable}
}

\begin{center}\rule{0.5\linewidth}{0.5pt}\end{center}

\subsection{Canonical 문서}\label{canonical-uxbb38uxc11c}

Protocol semantics의 기준은 \texttt{Protocol Spec.tex}이고,
구현 세부사항과 Open Decisions의 기준은
\texttt{Dev Spec.tex}이다.

\section{DPP Protocol Decision Log -
상세}\label{dpp-protocol-decision-log---uxc0c1uxc138}

이 장은 Recall에 한정하지 않고, 기존 GPT 대화와 이후 논의에서 나온 DPP protocol
전체 의사결정 과정을 기록한다. 이 문서는 ``왜 그렇게 결정했는가''를 보존하며,
현재 동작의 기준은 Protocol Specification과 Development Specification이다.

\begin{center}\rule{0.5\linewidth}{0.5pt}\end{center}

\subsection{1. 원본 대화와 문서
범위}\label{1-uxc6d0uxbcf8-uxb300uxd654uxc640-uxbb38uxc11c-uxbc94uxc704}

기존 GPT와 나눈 긴 수식/설계 대화는 workspace 안에 자동 편입된 것이
아니라 Downloads에 있다.

원본 파일은 \texttt{/Users/bighim/Downloads/ChatGPT-DPP} 수식 정리
\texttt{.md}이다.

확인된 규모는 약 16,393 lines이고, 31번의 prompt/answer interaction으로
구성되어 있다. 이 파일에는 Recall만이 아니라 proof system, Sahai
positioning, Pedersen commitment, DPP state vector, event relation,
provenance/privacy tradeoff 등 전체 설계 고민이 들어 있다.

따라서 이번 문서 체계의 기준은 다음과 같다.

\begin{description}[leftmargin=34mm,style=nextline]
  \item[Decision Log] 사람의 의도, 고민 과정과 반려된 선택지를 기록한다.
  \item[Protocol Spec] 결정 결과로 만들어진 DPP protocol의 canonical summary다.
  \item[Development Spec] Protocol Spec을 실제 구현 가능한 기능 단위로 풀어쓴다.
\end{description}

\begin{center}\rule{0.5\linewidth}{0.5pt}\end{center}

\subsection{2. Proof System과 구현 Stack
결정}\label{2-proof-systemuxacfc-uxad6cuxd604-stack-uxacb0uxc815}

처음 논의에서는 Groth16, Plonk, HyperPlonk, recursion, aggregation,
public input cost, verifier gas cost를 비교했다.

결론은 기존 \texttt{carbon-footprint} POC와 연속성을 유지하는 것이다.

\begin{itemize}
  \item Go와 gnark v0.15.0
  \item PLONK, BLS12-381과 KZG
  \item Poseidon2
  \item Solidity와 Foundry
\end{itemize}

Groth16은 proof size와 verifier gas cost에서 강점이 있지만, 현재 구현
연속성, universal setup 성격, gnark 기반 prototype 경험을 고려해 Plonk를
기본으로 둔다.

Aggregation, recursion, folding, batch verification은 v1 main
contribution이 아니다. 이들은 proof count나 verifier cost를 줄이기 위한
optimization 또는 future work로 둔다.

핵심 contribution은 proof system 자체가 아니라 다음이다.

\begin{itemize}
  \item DPP-specific private state payload
  \item Entry, Transfer, Merge, Split, Process와 Exit의 event-specific consistency relation
  \item Recall-default Transfer resolution layer
\end{itemize}

\begin{center}\rule{0.5\linewidth}{0.5pt}\end{center}

\subsection{3. Sahai 논문에 대한 이해와
Positioning}\label{3-sahai-uxb17cuxbb38uxc5d0-uxb300uxd55c-uxc774uxd574uxc640-positioning}

Sahai의 supply-chain 모델은 document graph에 가깝다.

\[
\begin{aligned}
\mathit{Document} &:= (\mathit{Type},\mathit{InDocs},\mathit{DocInfo}),\\
\mathit{DocInfo} &:= (\mathit{Item},\mathit{Sender},\mathit{Recipient},
\mathit{Quantity},\ldots).
\end{aligned}
\]

\texttt{DocInfo}는 직접 공개하지 않고 hash 또는 commitment로 숨긴다.
\texttt{InDocs}와 \texttt{OutDocs}는 provenance 추적을 위해 ledger에
남는다.

Sahai가 주장하는 privacy는 product graph 자체를 숨기는 의미가 아니라,
document 내부 값과 participant identity를 숨기는 의미에 가깝다. 즉
provenance graph는 존재하지만, graph의 각 node 내부 정보는 숨기는
구조다.

Sahai의 한계로 판단한 부분은 다음이다.

\begin{enumerate}
  \item Document predicate proof 중심이므로 input document와 검증 attribute가 늘면
  proof가 커진다.
  \item Merge와 Process의 의미가 충분히 분리되어 있지 않다.
  \item Process를 단순 Merge처럼 다루면 product transformation을 표현하기 어렵다.
  \item DPP-specific attribute를 위한 state-vector relation이 명확하지 않다.
  \item Hyperledger Fabric 기반이므로 public-chain contract와 gas model에는 차이가 있다.
\end{enumerate}

우리 모델의 위치는 Sahai를 부정하는 것이 아니라 확장하는 것이다.

Sahai-style document graph를 DPP-specific private state payload와 event-specific
relation으로 구체화한다.

\begin{center}\rule{0.5\linewidth}{0.5pt}\end{center}

\subsection{4. Token Recipes와 Event 기반 Provenance에서 얻은
결정}\label{4-token-recipesuxc640-event-uxae30uxbc18-provenanceuxc5d0uxc11c-uxc5bbuxc740-uxacb0uxc815}

Token Recipes 계열은 product transition을 token recipe, merge, split,
process, consume/new-token 형태로 표현한다. 중요한 교훈은
ERC20/ERC721만으로는 supply-chain lot transition을 자연스럽게 표현하기
어렵다는 점이다.

우리 모델도 다음 입장을 채택한다.

Product unit은 단순 NFT 하나가 아니다. Quantity와 state attribute를 가진 product
note 또는 slot-like asset으로 표현한다.

Token Recipes는 enforcement에 필요한 최소 상태는 contract storage에
두고, provenance 복원에 필요한 구조는 event log를 통해 재구성할 수 있게
한다. 그러나 이 방식은 privacy-preserving protocol이라기보다는
structural provenance model에 가깝다.

이 점은 우리 protocol의 privacy claim을 정하는 데 영향을 주었다.

Event-based public provenance를 선택하면 graph와 spent time은 public하다. 따라서
privacy claim은 product graph hiding이 아니라 content, state와 witness hiding이다.

\begin{center}\rule{0.5\linewidth}{0.5pt}\end{center}

\subsection{5. Hash-only Document와 Attribute Linking
구분}\label{5-hash-only-documentuxc640-attribute-linking-uxad6cuxbd84}

초기 가설은 supply-chain crypto model을 두 축으로 나누는 것이었다.

\begin{description}[leftmargin=45mm,style=nextline]
  \item[Hash-Based Document] Document나 data를 hash로 binding하고 integrity, audit,
  redaction과 update를 다룬다.
  \item[Attribute Linking / Consistency Check] Product state attribute가 Transfer,
  Merge, Split과 Process를 거치며 올바른 relation을 만족하는지 검증한다.
\end{description}

이 구분은 strict taxonomy라기보다는 recurring design pattern으로 보는
것이 더 안전하다.

Hash-only 방식은 document integrity와 off-chain data binding에는
충분하지만, 마지막 수령자가 특정 속성이 provenance로부터 올바른 도출
수식을 통해 자신에게 전달되었는지를 직접 검증하기 어렵다.

우리 모델은 Attribute Linking / Consistency Check 축에 있다.

DocumentHash는 일반 product information을 binding한다. Quantity와 StateVector는
$cm$ 안에 private witness로 binding되고 event circuit의 relation 검증에 사용된다.

\begin{center}\rule{0.5\linewidth}{0.5pt}\end{center}

\subsection{6. DPP Data Model 결정}\label{6-dpp-data-model-uxacb0uxc815}

DPP product information을 두 층으로 나눈다.

\begin{description}[leftmargin=31mm,style=nextline]
  \item[DocumentInfo] 반복 검증하지 않는 일반 product information
  \item[quantity] Split, Transfer와 Merge relation의 기준이 되는 수량
  \item[StateVector] Quantity에 비례 배분되거나 event relation으로 검증되는
  consistency-critical amount attribute
\end{description}

현재 notation은 다음이다.

\[
\begin{aligned}
usk &:= sk_{own}, & pk_{own} &:= H(sk_{own}),\\
addr &:= H(pk_{own}), & upk &:= (addr,pk_{own}),\\[2pt]
\mathit{DocumentInfo} &:= (d_1,d_2,\ldots,d_n),\\
\mathit{DocumentInfoBytes} &:=
\operatorname{EncodeText}(d_1)\mathbin{\|}\cdots\mathbin{\|}\operatorname{EncodeText}(d_n),\\
\mathit{DocumentHash} &:= \HashToField(\mathit{DocumentInfoBytes}),\\[2pt]
q &:= \mathit{quantity}, & \mathbf{s} &:= (s_1,s_2,\ldots,s_k),\\
\mathit{Document} &:= (\mathit{DocumentHash},q,\mathbf{s}),\\
cm &:= \PoseidonTwo(\mathit{DocumentHash},q,s_1,\ldots,s_k,addr,o),\\
nf &:= H(sk_{own},cm).
\end{aligned}
\]

Battery POC의 예시 state는 다음이다.

\[
q\in\mathbb{Z}_{\geq 0},
\qquad
\mathbf{s}_{battery}:=(total_{kg},carbon_{kgCO2e},recycled_{kg}).
\]

POC에서 \texttt{DocumentInfo}는 다음 field로 시작한다.

\[
\begin{aligned}
\mathit{DocumentInfo} &:= (\mathit{ProductName},\mathit{LotId},\mathit{Unit}),\\
b_x &:= \operatorname{UTF8}(\operatorname{NFC}(x)),\\
\operatorname{EncodeText}(x) &:={}
\operatorname{U32BE}(|b_x|)\mathbin{\|}b_x,\\
\mathit{DocumentInfoBytes} &:={}
\operatorname{EncodeText}(\mathit{ProductName})
\mathbin{\|}\operatorname{EncodeText}(\mathit{LotId})
\mathbin{\|}\operatorname{EncodeText}(\mathit{Unit}),\\
\mathit{DocumentHash} &:= \mathrm{fr.Hash}(\mathit{DocumentInfoBytes},\varepsilon,1)[0].
\end{aligned}
\]

\texttt{DocumentHash}는 product name, lot id, unit 같은 일반 product
identity 정보를 묶는다. 각 문자열을 개별 field element로 변환하지 않고,
각 문자열을 NFC 정규화한 UTF-8 byte string으로 만들고, unsigned 32-bit
big-endian byte 길이와 UTF-8 bytes를 순서대로 연결한 뒤 그 bytes를
\texttt{HashToField}로 field element 하나에 binding한다. Layout version과 별도
hash 용도 구분값은 사용하지 않는다. \texttt{fr.Hash} API의 DST에는 빈 byte string
$\varepsilon$을 전달한다.
\texttt{quantity}와 \texttt{StateVector}는 public document field가
아니라 note opening에 포함되는 private witness다.

\texttt{cm}은 finalized product note commitment다. \texttt{nf}는
finalized note를 소비할 때 double-spend를 막기 위한 nullifier다.

StateVector에 넣을 값은 quantity에 따라 분배 가능한 amount-like
attributes 중심으로 둔다. 현재 battery POC 예시는 \texttt{total\_kg},
\texttt{carbon\_kgCO2e}, \texttt{recycled\_kg}다. \texttt{quantity}는
StateVector에서 분리해 별도 field로 둔다. 이 값들은
Merge/Split/Transfer에서 additive 또는 quantity-proportional relation을
정의하기 쉽다.

\texttt{StateVector}라는 표현은 논문/스펙의 일반화 notation이다.
구현에서는 Poseidon2 input을 nested vector로 넣는 것이 아니라 다음
순서로 flat하게 넣는다.

\[
cm:=\PoseidonTwo(
\mathit{DocumentHash},q,total_{kg},carbon_{kgCO2e},recycled_{kg},addr,o).
\]

여기서 input ordering을 고정한다는 말은 serialization과 circuit witness
ordering을 고정한다는 의미다. Go prover, gnark circuit, Solidity
verifier/public input preparation, off-chain note reconstruction이 모두
같은 순서로 값을 해석해야 같은 \texttt{cm}이 나온다.

\texttt{cm}에는 \texttt{productType}, \texttt{version},
\texttt{schemaId} 같은 추가 metadata를 넣지 않는다. Product를 식별하거나
versioning에 필요한 값은 \texttt{DocumentInfo}에 넣고
\texttt{DocumentHash}로 binding한다. \texttt{StateVector}에는
consistency-critical state만 둔다.

반대로 percentage, grade, certification status, manufacturer, process
type처럼 quantity에 따라 자연스럽게 분배되지 않는 값은 State에 직접 넣지
않는 방향이 더 적절하다. 필요한 경우 percentage는 mass amount로 변환해서
넣고, grade/certificate/policy 같은 값은 \texttt{DocumentInfo}, policy,
별도 predicate, 별도 commitment category로 분리한다.

POC address model은 단순하게 둔다.

\[
usk:=sk_{own},\qquad
pk_{own}:=H(sk_{own}),\qquad
addr:=H(pk_{own}),\qquad
upk:=(addr,pk_{own}).
\]

사용자가 최종 정정한 POC 결정은 $addr=H(pk_{own})$이다. 즉
$pk_{own}=H(sk_{own})$을 먼저 만들고, address는 다시
$pk_{own}$을 hash해서 만든다. 이렇게 하면 기존
$upk=(addr,pk_{own})$ 인터페이스를 유지하면서도 address와
spend public key의 역할을 분리한다.

이 결정은 privacy-preserving one-time address 구조가 아니다.
\texttt{upk}에 \texttt{addr}와 \texttt{pk\_own}을 함께 두면 둘의 대응
관계는 공개된다. 이 POC 결정은 implementation simplification이고, 최종
identity/privacy model은 별도 설계 항목으로 남긴다.

\begin{center}\rule{0.5\linewidth}{0.5pt}\end{center}

\subsection{7. Pedersen Commitment 제거
결정}\label{7-pedersen-commitment-uxc81cuxac70-uxacb0uxc815}

초기 설계에서는 \texttt{C\_state}를 Pedersen commitment vector로 두고,
Merge/Split에서 외부 homomorphic check를 사용하는 방향을 검토했다.
그러나 Split에서 비례분배를 검증하려면 circuit이 \texttt{quantity},
input \texttt{StateVector}, output \texttt{StateVector}를 witness로
알아야 한다. 이 witness가 Pedersen commitment에 실제로 묶인 값이라는
것도 circuit 내부에서 증명해야 한다.

즉 Pedersen을 쓰더라도 $C_{state}\leftrightarrow\mathbf{s}$ witness의
binding relation을 증명해야 한다.

이 link를 증명하려면 circuit 내부에서 Pedersen opening 또는 vector
commitment relation을 검증해야 하며, 이는 scalar multiplication/MSM
기반의 commit-and-prove overhead를 만든다. 외부 homomorphic check 하나를
위해 이 overhead를 감수하는 구조는 현재 POC 목표와 맞지 않는다고
판단했다.

따라서 POC 설계에서는 Pedersen \texttt{C\_state}를 제거한다.
\texttt{quantity}와 \texttt{StateVector}는 note payload의 private
witness로 두고, note commitment가 이 값을 직접 binding한다.

\[
\mathit{Document}:=(\mathit{DocumentHash},q,\mathbf{s}),
\qquad
cm:=\PoseidonTwo(\mathit{DocumentHash},q,s_1,\ldots,s_k,addr,o).
\]

이 방향에서는 Merge/Split/Process relation을 circuit 내부에서 직접
검증한다.

\paragraph{Merge}
\[
q_{out}=q_{in,1}+q_{in,2},
\qquad
\mathbf{s}_{out}=\mathbf{s}_{in,1}+\mathbf{s}_{in,2}.
\]

\paragraph{Split}
\[
q_{in}=q_{out,1}+q_{out,2},
\qquad
q_{out,j}s_{in,k}=q_{in}s_{out,j,k}.
\]

\paragraph{Process}
Fixed-arity ProcessCircuit 안에서 aggregate, process delta, StateAllocationMatrix와
output StateVector relation을 검증한다.

잃는 것은 외부에서 arbitrary input/output에 대한 Pedersen homomorphic
aggregate check를 수행하는 능력이다. 하지만 현재 방향에서는 Split은
1-to-2 반복, Merge는 2-to-1 반복, Process는 max input/output fixed-arity
circuit으로 처리하므로 이 손실을 수용한다.

Proof는 다음을 함께 보장해야 한다.

\begin{itemize}
  \item $cm$ opening과 note binding
  \item ownership과 nullifier correctness
  \item Merkle membership과 range
  \item proportionality와 event-specific relation
\end{itemize}

현재 열린 질문은 \texttt{StateVector} encoding과 Poseidon2 note
commitment layout이다.

\begin{enumerate}
  \item Quantity와 StateVector를 field element로 어떻게 encode할 것인가?
  \item Poseidon2 arity와 input ordering을 어떻게 고정할 것인가?
  \item Attribute별 range를 어떻게 정할 것인가?
\end{enumerate}

이 질문은 아직 확정하지 않고
\texttt{Dev Spec.tex}의 Open Decisions에 둔다.

\subsubsection{Fixed-point와 decimal attribute
결정}\label{fixed-pointuxc640-decimal-attribute-uxacb0uxc815}

ZK circuit과 note commitment는 floating point 값을 직접 다루지 않는다.
POC에서는 \texttt{quantity}는 integer count로 두고, kg, kgCO2e처럼
소수점이 필요한 physical attributes는 nano scale fixed-point integer로
encode한다.

예를 들어 $q=3$은 정수 3으로 encode한다. 질량
$1.234567891\,\mathrm{kg}$은 scale $10^9$를 적용해 정수
$1{,}234{,}567{,}891$로 encode한다.

같은 attribute에 같은 scale을 일관되게 사용하고 field modulus를 넘지
않도록 range를 제한하면, fixed-point integer에 대한
addition/proportionality relation을 circuit 안에서 정확히 검증할 수
있다.

\[
\operatorname{Encode}(a+b)=\operatorname{Encode}(a)+\operatorname{Encode}(b).
\]

다만 Split에서 99:1 같은 비율 분배를 할 때 fixed-point 단위에서 정확히
나누어떨어지지 않을 수 있다. 처음에는 exact divisibility만 허용하고 정수 해가
없으면 proof 생성을 실패시키기로 했다. 이후 remainder를 첫 output에 할당하는
방식과 버리는 방식을 비교했다. Sound한 floor correctness를 유지하면서 검사를
줄이고 aggregate loss를 없애기 위해 deterministic first-output remainder
assignment를 최종 선택했다.

현재 확정된 범위는 다음뿐이다.

\[
0\leq q<2^{64},
\qquad
0\leq s_k<2^{64},
\qquad
\mathrm{Scale}_{State}=10^9.
\]

$q$는 scale 없는 integer count이고, StateVector는 nano scale fixed-point
integer다. Transfer의 첫 output은 $rv$, 두 번째 output은 $cm_{change}$이고
$rv$가 residual을 받는다. Split에서는 $cm_{out,1}$이 residual을 받고
$cm_{out,2}$를 내림 계산한다. Pedersen 외부 homomorphic check는 사용하지 않고
circuit 내부 floor/remainder relation과 aggregate conservation으로 검증한다.

Split output의 quantity가 0인지 별도로 검사하지 않는다. $q_1>0$ 또는 $q_2>0$은
POC 핵심 consistency relation에 필요하지 않은 semantic policy이므로 추가하지 않는다.
각 output에는 unsigned 64-bit range와 $q_{in}=q_1+q_2$를 적용한다.

Entry에서는 $0<q<2^{64}$와 $0\leq s_k<2^{64}$를 검사한다. Entry quantity는
0을 허용하지 않지만 각 $s_k=0$은 허용하며, attribute 간 별도 initial product
semantic invariant는 추가하지 않는다. 이 결정은 Split과 Transfer의 remainder
policy와 독립적이다.

\begin{center}\rule{0.5\linewidth}{0.5pt}\end{center}

\subsection{8. DPP Event Semantics
결정}\label{8-dpp-event-semantics-uxacb0uxc815}

기본 DPP event는 Entry, Transfer, Merge, Split, Process와 Exit의 6개다.

Recall과 Proceed는 별도 supply-chain transformation이 아니라, Transfer가
만든 pending state를 처리하는 resolution event다.

\subsubsection{Entry}\label{entry}

Entry는 외부 세계의 product information을 protocol 안의 finalized note로
만든다.

\[
\varnothing\longrightarrow cm_{new}.
\]

Entry의 quantity는 $0<q<2^{64}$이고 State coordinate는
$0\leq s_k<2^{64}$이다. 동일한 DocumentHash를 다시 사용하는 것은 허용한다.
각 note는 fresh opening $o$를 사용하므로 동일한 DocumentHash에서도 서로 다른 $cm$을
만들 수 있다. 중복 금지 대상은 DocumentHash가 아니라 정확히 같은 $cm$이다.

\subsubsection{Transfer}\label{transfer}

최신 결정에서 Transfer는 recall-default다. Transfer 시점에 receiver용
finalized \texttt{cm}을 만들지 않는다. 대신 receiver에게 보내는
portion은 \texttt{rv\_out}으로 만들고, sender에게 남는 portion은
\texttt{cm\_change}로 만든다.

\[
cm_{old}\longrightarrow (rv,cm_C).
\]

Transfer 표기에서 $T$는 receiver에게 이전하는 pending portion, $C$는 sender change,
$S$와 $R$은 Sender와 Receiver, $k$는 State coordinate index다. $r_{C,k}$는
$q_Cs_{in,k}$를 $q_{in}$으로 나눈 Euclidean remainder이며 Circuit witness로만 사용한다.
Note, voucher와 transaction에는 저장하지 않는다.

생성 Event에서 검증된 객체만 $MT$와 $rvMT$에 append한다는 closed-state invariant에
따라 Transfer는 input $q_{in}$과 $s_{in,k}$의 unsigned range를 반복 검사하지 않는다.
Output인 $q_T,q_C,s_{T,k},s_{C,k}$의 range는 검사한다. $q_T>0$은 zero transfer를
금지하기 위해 유지한다. $q_{in}>0$은 quantity conservation과 $q_T>0$에서 따라오고,
$q_C\geq0$은 unsigned range에서 따라오므로 별도 assert를 두지 않는다.

Zcash식 blind transfer처럼 동일한 note commitment를 계속 유지하는 방향은
현재 모델과 맞지 않는다. 이 protocol은 event-based public provenance를
선택했기 때문에 commitment-level product graph와 spent time은 public한
것으로 본다. 따라서 Transfer에서 sender에게 남는 portion을 숨기기 위해
같은 \texttt{cm}을 유지하기보다, explicit change note인
\texttt{cm\_change}를 생성하는 구현 방향을 택한다.

\subsubsection{Merge}\label{merge}

동일한 \texttt{DocumentHash}를 가진 product fragments만 merge한다.

\[
(cm_{old,1},cm_{old,2})\longrightarrow cm_{new},
\qquad
\mathbf{s}_{new}=\mathbf{s}_{old,1}+\mathbf{s}_{old,2}.
\]

\subsubsection{Split}\label{split}

하나의 product note를 여러 output note로 나눈다. \texttt{DocumentHash}는
유지한다.

\[
cm_{old}\longrightarrow(cm_{new,1},cm_{new,2}).
\]
\[
q_{new,1}+q_{new,2}=q_{old},
\qquad
q_{new,2}s_{old,k}=q_{old}s_{new,2,k}+r_{2,k}.
\]
\[
0\leq r_{2,k}<q_{old},
\qquad
s_{new,1,k}+s_{new,2,k}=s_{old,k}.
\]

이는 처음 검토했던 exact-divisibility 식을 대체한 최신 relation이다. 두 번째 output을
내림 계산하고 첫 번째 output이 residual을 받아 aggregate State 합을 보존한다.

\subsubsection{Process}\label{process}

Process는 input products가 output products로 변환되는 event다.
\texttt{DocumentHash}, product identity, quantity가 바뀔 수 있다.
Input과 output은 서로 다른 material/product일 수 있으므로 Process에서는 DocumentHash
equality나 transition relation을 검사하지 않는다. 각 DocumentHash가 해당 commitment에
binding되는지만 검사한다. Quantity transition도 POC에서 강제하지 않는다.

사용자가 재정리한 최신 Process relation은 input State를 output에 바로 분배하는
구조가 아니다. Input StateVector를 aggregate하고 공정 변화분을 반영한 뒤, 그 결과를
output product들에 allocation한다.

Notation은 다음으로 둔다. Input index는 $i$, output index는 $j$, State coordinate는
$k$, StateVector 길이는 $\ell$이다.

\[
\begin{aligned}
\mathbf{s}_{in}&:=\sum_{i=1}^{m}\mathbf{s}_{in,i},\\
\mathbf{s}_{out}&:=\sum_{j=1}^{n}\mathbf{s}_{out,j},\\
\mathbf{P} &\in \mathbb{Z}^{\ell},\\
A &\in \mathbb{Z}^{n\times\ell},\\
s_{out,k}A_{j,k} &= Ds_{out,j,k}+r_{j,k},
&& 2\leq j\leq n,\;1\leq k\leq\ell,\\
0\leq r_{j,k}&<D,
&& 2\leq j\leq n,\;1\leq k\leq\ell,\\
s_{out,1,k}+\sum_{j=2}^{n}s_{out,j,k}&=s_{out,k},
&& 1\leq k\leq\ell.
\end{aligned}
\]

수학적으로 하나의 signed vector $\mathbf{P}$만으로 coordinate별 증가, 감소와 보존을
모두 표현할 수 있다. 예를 들어 $(-10,+5,0)$은 첫 State를 10 감소시키고,
두 번째 State를 5 증가시키며, 세 번째 State를 보존한다. POC에서는 generic signed
Field encoding을 사용하지 않는다. State별 방향을 고정하고

\[
P=(-p_{totalLoss},+p_{carbonAdded},+p_{recycledDelta}),
\qquad
P_{pub}=(p_{totalLoss},p_{carbonAdded},p_{recycledDelta})
\]

로 encode한다. 첫 두 magnitude는 각각 66-bit range를 검사한다.
$M_{\max}=N_{\max}=3$이고 State coordinate가 64-bit이므로 aggregate의 상한은
$3(2^{64}-1)<2^{66}$이다. 66-bit 검사는 Field modular wraparound를 배제한다.
$p_{recycledDelta}$는 public input으로 전달하지만 별도로 0인지 검사하지 않는다. 세 번째 값을
생략하지 않는 이유는 StateVector coordinate와 public delta 배열의 위치를 일대일로
대응시키기 위해서다. Circuit은 다음 덧셈 relation을 검증한다.

Protocol Relation에서는 이를 한 줄의 vector equality로 나타낸다.

\[
\mathbf{s}_{out}=\mathbf{s}_{in}+P
\]

Development Specification은 public unsigned magnitude의 실제 circuit 구현을 명확히 하기
위해 같은 조건을 다음 coordinate별 식으로 전개한다.

\[
\begin{aligned}
s_{in,total}&=s_{out,total}+p_{totalLoss},\\
s_{out,carbon}&=s_{in,carbon}+p_{carbonAdded},\\
s_{out,recycled}&=s_{in,recycled}+p_{recycledDelta},\\
s_{out,recycled}&=s_{in,recycled}.
\end{aligned}
\]

마지막 두 relation을 함께 만족하면 $p_{recycledDelta}=0$이 도출된다. 따라서 같은
조건을 직접 반복하는 별도 0 assertion은 두지 않기로 결정했다.

여기서 $\mathbf{s}_{in,i}$와 $\mathbf{s}_{out,j}$는 각 note의 StateVector이고,
$\mathbf{s}_{in}$과 $\mathbf{s}_{out}$은 각각의 aggregate StateVector다.
별도 $\mathbf{S}_{in},\mathbf{S}_{out}$ 행렬은 이후 수식에 사용되지 않아 제거했다.
$\mathbf{P}$는 process delta vector, $A$는 StateAllocationMatrix다.
$D$는 state value scale이 아니라 allocation ratio를 정수로 표현하기
위한 fixed-point allocation denominator이며, POC에서는 constant $D=10^9$으로 고정한다.
$A$는 public input으로 proof와 transaction에 포함한다.

Process는 output 2부터 $n$까지 내림 quotient와 remainder를 검증하고,
output 1이 모든 State coordinate의 residual을 받는다. 따라서 aggregate State는
버려지지 않는다. State coordinate마다 $n-1$개의 floor relation과 remainder
range check, 한 개의 aggregate conservation equality를 사용한다. POC maximum인
$n=3$, $\ell=3$에서는 floor relation과 remainder range check가 각각 6개,
conservation equality가 3개다.

Allocation matrix는 attribute별로 output 배분 방식을 다르게 표현할 수
있어야 한다. 예를 들어 product D와 waste E가 나오는 process에서 carbon과
recycled는 product D에 전부 배정하고, total\_kg만 D/E에 비례 배분하는
식의 recipe를 표현할 수 있다.

Process circuit은 public \texttt{P\_pub}를 입력받는다. Circuit은 input과 output에서
각각 $\mathbf{s}_{in}$과 $\mathbf{s}_{out}$을 계산하고
$\mathbf{s}_{out}=\mathbf{s}_{in}+\mathbf{P}$에 대응하는 coordinate별 relation을
확인한 뒤, $\mathbf{s}_{out}$이 $A$에 따라 각 $\mathbf{s}_{out,j}$로 정확히
배분되었는지를 확인한다. 예를 들어 battery POC에서는 다음 종류의
aggregate constraint를 둘 수 있다.

\[
s_{out,carbon}\geq s_{in,carbon},
\qquad
s_{out,recycled}=s_{in,recycled},
\qquad
s_{out,total}\leq s_{in,total}.
\]

위 constraint는 output aggregate $\mathbf{s}_{out}$에 직접 적용된다. Active Process
output에는 $1\leq q_{out,j}<2^{64}$를 요구한다. Quantity가 0인 output을 만들지 않고,
필요한 output 수만 $n$으로 전달하기로 결정했다.

Recipe는 POC에서는 \texttt{StateAllocationMatrix\ A} 자체로 단순화할 수
있다. 더 일반적인 논문 표현에서는 \texttt{A}를 process recipe의
state-allocation component로 둘 수 있다. POC에는 on-chain recipe registry나
승인된 recipe 목록을 두지 않고, $A$를 public input으로 공개한다. 따라서 실제 사용된
allocation은 숨길 수 없지만, 호출자가 선택한 $A$ 자체의 물리적 타당성은 보장하지 않는다.

이 POC가 막는 조작과 막지 못하는 조작을 구분한다.

\begin{itemize}
  \item Input commitment에 binding된 State를 임의로 바꾸거나 이미 소비한 note를 다시
  사용하는 것은 막는다.
  \item Output aggregate에서 carbon 감소, recycled 총량 변경과 total 증가를 막는다.
  \item Output State가 transaction에 공개된 $A$와 다르게 배분되는 것을 막는다.
  \item Entry 또는 실제 공정에서 허위 State가 입력되는 것은 막지 못한다.
  \item 공개된 $P$와 다른 carbon 증가량을 output에 반영하는 것은 막는다.
  다만 실제 공정에서 발생한 carbon보다 작은 $P$를 허위로 공개하는 것은 외부
  measurement나 attestation이 없으면 막지 못한다.
  \item 호출자가 carbon 또는 recycled를 특정 output에 집중시키는 $A$를 선택하는 것은
  막지 못한다. Public $A$는 이를 탐지 가능하게 만들 뿐 금지하지 않는다.
  \item 특정 output에 State를 몰아준 뒤 Exit하는 행위를 막으려면 $A$에 대한 별도
  물리 정책, output type binding 또는 Exit policy가 필요하며 현재 POC 범위에는 없다.
\end{itemize}

\[
Recipe_R\Longrightarrow A_R.
\]

즉 recipe가 allowed input/output type, optional quantity policy, process
delta policy, state allocation rule을 포함하고, 그 중 state allocation
component가 \texttt{A\_R}가 되는 구조다. POC에서는 이 중 \texttt{A}
중심으로 먼저 구현한다.

Process는 inverse allocation 관점으로 설계하지 않는다. \texttt{A}는 보통
square matrix가 아니고, 특정 attribute가 어떤 output에 전혀 배정되지
않는 경우 $A_{j,k}=0$이 가능하므로 inverse가
일반적으로 존재하지 않는다. 따라서 circuit은 output에서 intermediate를
역산하지 않고, \texttt{I}에서 output \texttt{Y\_j}를 forward
allocation하는 relation을 검증한다.

Process는 외부 vector commitment equation에 기대지 않는다.
\texttt{StateAllocationMatrix\ A}와 coordinate-wise allocation
relation은 단순 additive homomorphism으로 검증할 수 없고, 결국 circuit
내부에서 state witness와 output commitment binding을 함께 확인해야 한다.
따라서 Pedersen vector commitment를 추가하면 circuit 내부 commitment
opening/MSM overhead만 커진다. 현재 POC 방향은 Pedersen을 제거하고
fixed-arity ProcessCircuit 내부에서 aggregate, process delta,
allocation, output commitment를 직접 검증하는 것이다.

\subsubsection{Exit}\label{exit}

Exit는 finalized note를 protocol 밖으로 내보내거나 폐기한다.

\[
cm_{old}\longrightarrow\varnothing.
\]

\begin{center}\rule{0.5\linewidth}{0.5pt}\end{center}

\subsection{9. Recall 관련
결정}\label{9-recall-uxad00uxb828-uxacb0uxc815}

Recall은 main contribution이 아니라 Transfer에 붙는 pending-transfer
resolution layer다.

최종 결정은 다음이다.

\begin{enumerate}
  \item $cm$ 안에는 RecallVoucher를 넣지 않는다.
  \item Transfer는 recall-default다.
  \item Transfer 시점에 $cm_{pending}$을 만들지 않는다.
  \item Transfer는 pending $rv$와 sender의 change note $cm_C$를 만든다.
  \item Proceed는 $rv$와 $rvnf$를 사용해 receiver용 $cm_{recv}$를 새로 만든다.
  \item Recall은 $rv$와 $rvnf$를 사용해 sender용 $cm_{return}$을 새로 만든다.
\end{enumerate}

이전에는 \texttt{cm\_pending}을 만들고 receiver가 그
\texttt{cm\_pending}을 소비하는 구조를 검토했다. 하지만 event-based
public provenance를 선택하면 graph와 spent time은 어차피 public이다.
또한 \texttt{cm\_pending}은 finalized note처럼 쓰이지 않고 pending
record 역할만 하므로 별도 note commitment로 만들 필요가 약하다.

따라서 pending transfer는 \texttt{rv}로 표현한다.

\[
\begin{aligned}
rv &:= H(\mathit{DocumentHash},q_T,\mathbf{s}_T,addr_{Sender},addr_{Receiver},
\Delta_{epoch},o_{rv}),\\
rvnf &:= H(rv,o_{rv}).
\end{aligned}
\]

\texttt{o\_rv}는 sender와 receiver만 아는 Recall Voucher opening
secret이다. 따라서 sender와 receiver만 같은 \texttt{rvnf}를 만들 수
있다.

\begin{center}\rule{0.5\linewidth}{0.5pt}\end{center}

\subsection{10. On-chain Storage와 Timing
결정}\label{10-on-chain-storageuxc640-timing-uxacb0uxc815}

pending transfer 상태는 enum status가 아니라 deadline mapping으로
관리한다.

\[
\Map{rv}[rv]=deadline.
\]

Transfer가 성공하면 $\Map{rv}[rv]$가 생성된다. Recall 또는 Proceed가
성공해도 이 값을 삭제하거나 덮어쓰지 않는다.

Recall 또는 Proceed가 성공하면 $rvnf$를 $\List{rvnf}$에 기록한다. Solidity에서는
동일한 의미를 bool mapping으로 구현할 수 있다.

검토 과정에서 raw voucher hash를 직접 key로 사용하는 $resolved[rv]$ bool mapping으로
단순화하는 대안을 비교했다. Proceed와 Recall이 public input으로 $rv$를 이미 전달하므로
$rvnf$가 추가적인 on-chain unlinkability를 제공하지 않는다는 지적은 맞다. 그럼에도
finalized note의 $nf$와 일관된 secret-derived nullifier 구조를 유지하고, $o_{rv}$를
공유한 Sender와 Receiver가 동일한 resolution identifier를 계산하도록 $rvnf$를
유지하기로 결정했다. 따라서 별도 $resolved[rv]$ mapping은 사용하지 않는다.

초기에는 \texttt{Pending}, \texttt{Recalled}, \texttt{Finalized} status
enum을 고려했다. 그러나 $\Map{rv}[rv]$는 deadline history를 보존하고,
resolution 중복 방지는 $\List{rvnf}$ membership으로 처리하는 것이 더 단순하다고
판단했다.

timing은 block timestamp 경계조건을 피하기 위해 epoch 기반으로 둔다.

\[
\begin{aligned}
EPOCH\_SIZE &:= 600\ \mathrm{seconds},\\
epoch_{transfer} &:=
\left\lfloor\frac{block.timestamp}{EPOCH\_SIZE}\right\rfloor,\\
deadline &:= epoch_{transfer}+\Delta_{epoch}.
\end{aligned}
\]

Recall은 $current_{epoch}<deadline$인 동안만 가능하고, Proceed는 Recall/Proceed로
resolve되지 않았다면 deadline 이후에도 가능하다. $\Delta_{epoch}=0$이면
$deadline=epoch_{transfer}$이므로 Recall은 처음부터 불가능하다. $\Delta_{epoch}$은
voucher hash에 binding되는 public Field Element지만 Circuit에서 별도 bit range check를
하지 않는다. Generated verifier가 $0\leq\Delta_{epoch}<r$인 canonical public input인지
검사한다. Main Contract는 같은 범위 검사를 반복하지 않고 proof 검증 후 deadline을
계산한다. $\Delta_{epoch}<r<2^{255}$이고 $currentEpoch<2^{247}$이므로 합은
\texttt{uint256} 범위 안에 있다.

\begin{center}\rule{0.5\linewidth}{0.5pt}\end{center}

\subsection{11. Provenance와 Privacy
결정}\label{11-provenanceuxc640-privacy-uxacb0uxc815}

이 절은 event-based public provenance를 검토하며 도출한 privacy 경계를 보존한다.
현재 POC에서는 event schema가 확정되지 않아 provenance event emission을 보류했지만,
향후 commitment edge를 공개하면 아래 분석이 그대로 적용된다.

따라서 commitment-level product graph, event execution time, receiver spent time과
직간접 downstream spent time은 숨기지 않는다.

이 결정은 Token Recipes와 Sahai를 검토한 결과다. provenance를 추적
가능하게 만들면 input/output commitment edge가 공개되고, 이 edge를
따라가면 downstream product graph도 복원 가능하다.

따라서 우리 논문에서 privacy claim은 다음으로 제한해야 한다.

\begin{description}[leftmargin=28mm,style=nextline]
  \item[숨기는 정보] DocumentInfo와 State의 구체 값, proof witness, 그리고 최종
  identity model이 지원하는 범위의 participant identity
  \item[숨기지 않는 정보] Commitment-level provenance graph, event timing과 spent timing
\end{description}

POC에서는 address를 숨기지 않는다. $pk_{own}=H(sk_{own})$,
$addr=H(pk_{own})$ 구조를 사용한다. 다만 최종
논문/프로토콜에서 participant identity와 protocol address를 얼마나
분리할지는 아직 확정하지 않는다.

\begin{center}\rule{0.5\linewidth}{0.5pt}\end{center}

\subsection{12. 현재 남은 큰 미해결
질문}\label{12-uxd604uxc7ac-uxb0a8uxc740-uxd070-uxbbf8uxd574uxacb0-uxc9c8uxbb38}

Pedersen \texttt{C\_state}를 제거하고 Process를 fixed-arity circuit
중심으로 재정리하면서, 남은 질문도 다음처럼 바뀐다.

\begin{enumerate}
  \item Event별 aggregate와 multiplication의 intermediate range
  \item Event별 circuit의 실제 intermediate bound를 test vector와 함께 확정
\end{enumerate}

여기에 구현 관점에서 다음 질문도 남아 있다.

\begin{itemize}
  \item 최종 Address/identity privacy model
  \item event schema와 indexed field
\end{itemize}

미확정 내용은 임의로 채우지 않고
\texttt{Dev Spec.tex}의 Open Decisions에 둔다.

\begin{center}\rule{0.5\linewidth}{0.5pt}\end{center}

\subsection{13. 후속 결정과 현재 유효 상태}

앞 절은 설계가 발전한 순서를 보존한다. 이후 논의에서 다음 결정이 기존 기록을
대체하거나 구체화했다.

\subsubsection{Curve와 On-chain Verifier}

초기 carbon-footprint POC의 BN254 연속성보다 BLS12-381을 사용하는 방향으로
개발 Stack을 변경했다. gnark v0.15.0의 실제 source와 최소 재현 실험을 통해
BLS12-381 PLONK proof 생성, Go verification, Solidity verifier export와
proof serialization을 확인했다. 이어서 Foundry v1.7.1의 Prague EVM에서
EIP-2537 G1 MSM과 pairing precompile을 호출하는 verifier가 valid proof를
수락하고 변조 public input을 거부하는 것을 확인했다.

따라서 현재 POC Stack은 다음으로 확정한다.

현재 POC Stack은 Go, gnark v0.15.0, PLONK, BLS12-381, KZG, Poseidon2,
Solidity와 Foundry Prague EVM으로 확정한다.

POC KZG SRS는 \texttt{unsafekzg.NewSRS}로 만들며 production setup으로 사용하지 않는다.

\subsubsection{Poseidon2 Parameter}

Poseidon2는 Field와 width, round 수, S-box, round constants와 hash construction을
선택해야 하는 parameterized hash family다. POC가 별도 instance를 설계하거나 자체
round constants를 만들지 않고, version-pinned gnark implementation의 default를
재사용하기로 결정했다.

\[
\begin{aligned}
\mathrm{Field} &= \mathrm{BLS12\text{-}381\ ScalarField},\\
t &= 2,\\
R_F &= 6,\\
R_P &= 50,\\
\alpha &= 5,\\
\mathrm{Construction} &= \mathrm{Merkle\text{-}Damg\mathring{a}rd}.
\end{aligned}
\]

대상 구현은 gnark v0.15.0과 그 dependency인 gnark-crypto v0.20.1이다. $H$와
$CM$의 Field Element input은 각 protocol 식의 순서로 전달하고, 별도 문자열 domain
tag는 넣지 않는다. Go와 circuit의 동일성을 확인하는 reference vector를 test artifact로
남긴다. Library version 또는 default parameter가 바뀌면 circuit과 PLONK key를 다시
생성해야 한다.

\subsubsection{Event Arity와 Nullifier 표기}

Merge와 Split을 arbitrary arity로 구현하지 않는다. Merge는 fixed 2-to-1,
Split은 fixed 1-to-2 circuit으로 구현하고 필요한 경우 반복 composition한다.
Finalized note nullifier와 Recall Voucher nullifier는 다음 표기를 사용한다.

\[
nf=H(sk_{own},cm),
\qquad
rvnf=H(rv,o_{rv})
\]

별도의 Poseidon2 해시 용도 구분값은 현재 POC 식에 넣지 않는다. On-chain protocol object의
존재 검사는 $\List{cm}$, $\List{nf}$, $\List{rv}$, $\List{rvnf}$와 같은
logical set notation으로 표현하고, 실제 Solidity representation은 Development
Specification에서 결정한다.

\subsubsection{Provenance Event와 Tree Index Event의 분리}

초기에는 public event edge로 provenance graph를 복원하는 방향을 채택했다. 그러나
event schema와 공개 범위를 아직 확정하지 않았으므로 현재 POC의 Smart Contract에서는
input/output commitment edge를 직접 연결하는 provenance event를 제외한다. 향후 이
edge를 공개하기로 결정하면 commitment-level graph와 spent time이 공개된다는 분석은
여전히 유효하다.

이 결정과 Merkle witness 획득용 event는 구분한다. Client가 자신이 받은 $cm$ 또는
$rv$의 leaf index를 찾을 수 있도록 Tree append 시 leaf, zero-based index와 new root를
event로 공개한다. 이 event만으로 old input과 new output 사이의 provenance edge를
직접 선언하지는 않는다. 현재 privacy claim은 DocumentInfo, StateVector와 witness
hiding 범위로 제한한다.

\subsubsection{문서 통합}

중복 문서의 수동 동기화를 줄이기 위해 canonical 문서를 다음 세 개로 단순화한다.

\begin{enumerate}
  \item \texttt{Protocol Spec.tex}: protocol semantics와 formal Scheme
  \item \texttt{Dev Spec.tex}: 구현 선택, ABI, artifact와 test
  \item \texttt{Decision Log.tex}: Summary와 의사결정 과정
\end{enumerate}

Paper Scheme과 HTML은 canonical source가 아니라 위 문서에서 파생하는 view로 취급한다.

\begin{center}\rule{0.5\linewidth}{0.5pt}\end{center}

\subsection{14. Specification 대조에서 보완한 결정}

Protocol Specification과 Development Specification을 현재 Decision Log와 대조한 결과,
명세에는 반영되어 있으나 결정 과정이 누락된 항목을 다음과 같이 보완한다.

\subsubsection{DocumentHash의 검증 경계}

DocumentInfo의 각 문자열을 Unicode NFC로 정규화하고 UTF-8로 encode한다. UTF-8
byte 길이를 unsigned 32-bit big-endian으로 먼저 기록한 뒤 해당 bytes를 이어 붙인다.
세 field는 ProductName, LotId, Unit 순서로 연결한다. 그 결과에 BLS12-381 scalar
field의 \texttt{fr.Hash}를 적용해 DocumentHash 하나를 만든다. \texttt{fr.Hash}는
SHA-256 기반 XMD를 사용하는 Hash-to-Field API이며 단순한
\texttt{SHA256(msg) mod r} 계산이 아니다. Layout version과 별도 hash 용도 구분값은
사용하지 않고, API의 DST에는 빈 byte string을 전달한다.

이 encoding은 다음 선택지를 비교해 결정했다.

\begin{itemize}
  \item 단순 concatenation은 \texttt{("ab","c")}와 \texttt{("a","bc")}의
        경계를 구분하지 못하므로 반려했다.
  \item delimiter 방식은 delimiter escaping과 canonical 처리 규칙이 추가되므로
        반려했다.
  \item variable-length integer는 최소 encoding 여부를 추가로 정의해야 하므로 POC에
        불필요한 복잡성을 만든다.
  \item fixed 4-byte length는 Go와 다른 언어에서 구현이 단순하고 POC field 크기에
        충분하다. Big-endian은 network byte order와 Solidity ABI의 정수 표현 방향에
        맞추기 위해 선택했다.
  \item normalization을 하지 않으면 시각적으로 같은 문자열의 NFC/NFD 표현이 서로
        다른 hash를 만든다. NFC는 이를 통일하면서 product code의 compatibility 문자를
        바꿀 수 있는 NFKC보다 의미 보존적이므로 선택했다.
\end{itemize}

길이는 Unicode code point 또는 rune 수가 아니라 NFC 적용 후 UTF-8 byte 수다. Invalid
UTF-8과 $2^{32}-1$ byte를 초과하는 field는 거부한다. Empty string 허용 여부는 encoding
문제가 아니라 DocumentInfo semantic validation의 문제다. 현재 POC는 별도 semantic
validation을 두지 않으므로 Empty string과 NUL byte가 포함된 유효한 UTF-8 field를 허용한다.

DocumentInfo는 circuit witness나 public input으로 넣지 않는다. Circuit은
DocumentHash가 input/output 사이에서 유지되거나 Process가 허용한 방식으로 변경되는지만
검증한다. Receiver는 전달받은 DocumentInfo를 직접 serialize하고 HashToField를 다시
계산해 DocumentHash와 일치하는지 off-chain에서 확인한다. 모든 leaf 정보를 contract에
올려 hash를 재계산하지 않는 것은 product information을 on-chain에 공개하지 않기 위한
명시적인 trust boundary다.

\subsubsection{Note, Owner Address와 Freshness}

Finalized note의 wallet-side representation은 다음과 같다.

\[
note=(cm,DocumentHash,q,\mathbf{s},o)
\]

Note에는 $addr$을 중복 저장하지 않는다. Holder는 KeyGen/Register 단계에서 정한
자신의 $addr$을 별도로 알고 있으며 Client 함수가 해당 값을 사용한다. Circuit은
$pk_{own}=H(sk_{own})$과 $addr=H(pk_{own})$ relation을 검증한다.

초기에는 Azeroth의 EOA mapping 패턴을 따라 $addr$ 중복을 거부하는 방식을 검토했다.
현재 POC는 이를 단순화해 \texttt{UserByEOA[msg.sender]}에 $(addr,pk_{own})$만 저장하고
uniqueness나 hash relation을 Contract에서 검사하지 않는다. 이 mapping은 wallet 조회를
위한 metadata이며, 공급망 Event submitter가 해당 EOA여야 한다는 조건으로 사용하지 않는다.

새 note마다 $sk_{own}$, $pk_{own}$과 $addr$을 새로 만들지 않는다. POC에서는 Client가
선택한 address를 재사용하되, 새 $cm$을 만들 때마다 fresh opening $o$를 생성한다. Commitment에
들어가는 address는 output의 실제 holder를 따른다.

\begin{itemize}
  \item Entry output: Entry가 지정한 이미 등록된 holder address
  \item Transfer change: $addr_{Sender}$
  \item Proceed output: $addr_{Receiver}$
  \item Recall output: Recall을 실행하는 직전 sender의 address
  \item Merge, Split, Process의 self-operation output: 실행 holder의 address
\end{itemize}

\subsubsection{Numeric Encoding과 Zero Change}

POC에서 $q$와 각 StateVector coordinate는 unsigned 64-bit integer로 제한한다.
$q$는 item count이고, $total_{kg}$, $carbon_{kgCO2e}$, $recycled_{kg}$는
nano scale $10^9$ fixed-point integer다. Aggregate와 cross multiplication의 더 큰
intermediate bound는 circuit별 range 분석 후 정한다.

Transfer는 $q_C=0$이어도 $cm_C$를 항상 생성한다. 이 경우 proportionality와 range
constraint에 의해 change StateVector도 zero vector여야 한다. 이 선택은 Transfer ABI와
output shape를 모든 입력에 대해 일정하게 유지한다.

\subsubsection{Process POC의 Maximum Arity}

Process의 일반 모델은 maximum input/output arity $M_{\max}$와 $N_{\max}$를 사용한다.
기능 확인용 기본 POC에서는 $M_{\max}=N_{\max}=3$으로 고정한다. 실제 active input $m$과 output $n$이
3보다 작으면 public active count $m,n$에 binding된 algebraic gating과 zero-padding
constraint로 나머지 slot을 비활성화한다. 별도 enable vector는 statement 또는 Client
witness에 넣지 않는다.

성능 측정에서는 같은 3-to-3 circuit에 inactive slot만 넣어 작은 arity의 비용으로 해석하지
않는다. $M_{\max}=N_{\max}=1$, 2, 3인 circuit을 각각 compile하고 기준 scenario에서
모든 slot을 활성화해 1-to-1, 2-to-2와 3-to-3을 측정한다. 비대칭 Process는 이번
benchmark에서 제외한다. 측정 결과를 확인하기 전에는 더 큰 maximum arity를 정하지 않는다.

\subsubsection{State 길이와 Process Benchmark Matrix}

StateVector constraint 증가량을 확인하기 위해 $\ell=0,1,2,3$인 별도 circuit profile을
compile한다. State layout은
$()$, $(total_{kg})$, $(total_{kg},carbon_{kgCO2e})$,
$(total_{kg},carbon_{kgCO2e},recycled_{kg})$ 순서의 prefix다. 같은 최대 circuit을 두고
unused State 값을 0으로 채우는 방식은 State constraint scaling을 측정하지 못하므로 사용하지
않는다.

State scaling은 Depth 32에서 수행하며 Process는 3-to-3으로 고정한다. Process arity
scaling은 State 3과 Depth 32에서 1-to-1, 2-to-2, 3-to-3을 측정한다.
따라서 $(\ell<3,a<3)$ 교차 profile은 생성하지 않는다. 이후 profile은 모두
Poseidon2를 사용한다.

공통 State-3/Depth-32/Process-3-to-3 결과는 State와 arity 비교표에서 재사용한다.
중복을 제거하면 각 Non-Process Event는 Poseidon2 State profile 4개이고,
Process는 Poseidon2 State profile 4개와 추가 arity profile 2개로 6개다.
전체 active 범위는 $7\times4+6=34$개 benchmark configuration이며 distinct circuit도
34개다. 각 configuration은 constraint 수,
proving/verifying time, artifact/proof size,
verifier gas와 protocol transaction gas를 포함하는 하나의 benchmark row다.

\subsubsection{두 Merkle Tree와 History 보존}

Finalized note membership 구조를 미확정으로 두었던 이전 기록은 다음 결정으로 대체한다.

\begin{itemize}
  \item $MT$: finalized $cm$을 append하는 공급망 note tree
  \item $rvMT$: pending $rv$를 append하는 Recall Voucher tree
  \item 두 Tree는 동일한 depth $d$를 사용한다.
  \item POC와 benchmark는 $d=32$로 고정하며 Depth scaling은 수행하지 않는다.
  \item Empty leaf는 0이며 중간 default node는 아래 default node 두 개를
        $H_{\mathrm{tree}}$로 hash해 재귀 계산한다.
  \item $cm$과 $rv$는 leaf-hash 없이 그대로 append한다. Valid output은 0이 아니지만
        Azeroth처럼 별도 nonzero 검사는 하지 않고 hash output이 0일 약 $1/r$의
        가능성을 무시한다.
  \item 두 Tree에서 생성된 모든 과거 root를 영구 허용한다.
\end{itemize}

Contract는 accepted note roots $\List{rt}$와 voucher roots $\List{rvrt}$를 구분한다.
Transfer는 $cm_C$를 $MT$에, $rv$를 $rvMT$에 각각 append한다. Proceed와 Recall은
$rvMT$ membership을 증명하고 $rvnf$를 기록해 같은 voucher의 두 번째 resolution을
거부한다.

$\Map{rv}[rv]=deadline$은 resolution 이후에도 삭제하지 않는다. $\List{rvnf}$도
history를 보존한다. Solidity에서는 set membership을 bool mapping으로 구현할 수 있지만,
formal Scheme은 자료구조 구현과 분리된 logical set notation을 유지한다.

Tree의 Contract layout은 Azeroth \texttt{BaseMerkleTree.sol} 구조를 따른다. Leaf를
왼쪽부터 순차 append하고 실제 leaf와 영향을 받은 모든 internal node를 heap-indexed
mapping에 저장한다. 기록되지 않은 0 storage slot은 해당 level의 재귀 default node로
해석한다. Internal hash는 Azeroth의 abstract \texttt{\_hash(left,right)} 경계를
유지하되 gnark-crypto v0.20.1 BLS12-381 default Poseidon2 direct width-2 compressor로
교체한다.

Membership witness는 다음처럼 정의한다.

\[
path=(index,sibling_0,\ldots,sibling_{d-1})
\]

Circuit은 $index$의 $d$-bit little-endian decomposition으로 각 level의 left/right
ordering을 선택한다. gnark가 이 selection constraint를 생성하더라도 Client는
$index$와 sibling $d$개를 제공해야 한다. gnark 기본 Merkle gadget은 leaf를 먼저
FieldHasher로 hash하므로 Azeroth식 raw leaf root와 일치하지 않는다. 따라서 circuit은
raw leaf에서 시작해 direct Poseidon2 compressor를 $d$번 적용하는 전용 membership
verifier를 사용한다.

Azeroth는 append 이후 증가한 \texttt{\_numLeaves}를 event의 \texttt{index}로
emit하여 실제 zero-based leaf index보다 1 큰 값을 내보낸다. zkDPP는 삽입 전 leaf
count를 보존하고 다음 event로 정확한 zero-based index를 공개한다.

\[
\mathrm{MTAppend}(cm,index,rt_{new}),
\qquad
\mathrm{RVMTAppend}(rv,index,rvrt_{new})
\]

\subsubsection{Hash Profile과 Event별 Benchmark}

POC와 benchmark의 depth는 32로 고정한다. 기능 확인용 기본 배포는 Poseidon2,
$\ell=3$, Process $a=3$인 8개 Relation을 사용한다. Active benchmark에서는 Poseidon2
State scaling 32개와 State-3/Depth-32 Process 1-to-1, 2-to-2 2개를 생성한다.
공통 기준점을 중복하지 않은 benchmark configuration과 distinct circuit 총수는 모두
34개다. SHA-to-Field profile은 active benchmark와 이후 POC 검토 범위에서 제외한다.

Depth scaling은 제거한다. 대신 depth-32 Merkle membership 하나만 독립적으로 측정해
commitment 한 개의 path verification에 드는 constraint, proving/verification time과
Solidity gas를 기록한다.

완료된 SHA profile의 single-Field layout은 다음과 같다. 각 입력 Field Element를 canonical
32-byte unsigned big-endian으로 encode하고 왼쪽에서 오른쪽 순서로 연결한다. SHA-256
digest를 unsigned big-endian integer로 해석한 뒤 BLS12-381 scalar modulus $r$로 환원한다.
따라서 SHA profile도 $cm$, $nf$, root와 public input 하나당 Field Element 한 개를
유지한다. Digest를 두 limb로 보존하는 방식은 public input과 Contract ABI 자체가 달라져
동일 조건의 hash scaling 비교가 아니므로 채택하지 않았다. 이 방식은 DocumentHash에
사용하는 \texttt{fr.Hash}와 다르며, modulo reduction의 분포 편향을 허용한 POC benchmark
선택이다.

SHA profile을 추가했던 주된 목적은 prover 성능 개선이 아니라 EVM gas 비교였다. Solidity의
SHA-256 precompile을 사용하는 compressor와 generated Poseidon2 Solidity 연산을 같은
depth-32 Tree에서 비교한다. 따라서 결과표는 circuit constraint와 prover 비용을 숨기지
않되, 핵심 on-chain 지표로 compressor 호출, hasher/Main Contract 배포, depth-32 append,
Entry verifier와 실제 Entry transaction gas를 분리해 기록한다.

각 Event에 대해 circuit constraint 수, compile/setup time, proving/verifying key
크기, proof 생성 시간, Go verification 시간, proof와 transaction 크기, verifier
deployment gas, verifier-only gas와 전체 protocol transaction gas를 측정한다.
Tree append gas와 witness/path 생성 시간도 별도 microbenchmark로 측정한다.

기준 의미론은 Entry $0\to1$, Transfer의 old note 한 개와 $cm_{change}/rv$, Proceed와
Recall의 voucher 한 개, Merge $2\to1$, Split $1\to2$, Process $a\to a$, Exit
$1\to0$으로 고정한다. Depth 비교는 State 3의 동일한 의미론적 input을 사용한다.

\subsubsection{첫 구현 Milestone의 범위와 결과}

전체 8개 Event를 동시에 구현하면 Poseidon2, Merkle Tree, generated verifier와 Main
Contract integration의 실패 원인을 분리하기 어렵다. 따라서 첫 milestone은
$\ell=3,d=32,m=n=3$인 Entry와 Process만 구현하기로 결정했다. Process statement에는
$m,n$을 유지하지만 circuit은 모두 3인지 직접 검사한다. selector와 inactive slot은
3-to-3 vertical slice가 통과한 뒤 추가한다.

기존 23-public-input Process로 Entry와 Process vertical slice 및 SHA-to-Field 비교를
시험했다. 이후 $p_{recycledDelta}$ coordinate를 포함하는 single-root 24-input 명세를
거쳤지만, 최종적으로 input별 $rt_i$를 포함하는 26-input 명세로 변경했다. 따라서 이전
23-input과 24-input 상태에서 얻은 constraint, time, size와 gas 수치는 canonical 결과에서
제거하고 재사용하지 않는다.

사용자는 SHA 비교를 종료하고 이후 POC 구현과 benchmark를 Poseidon2로만 진행하기로
결정했다. 새 benchmark는 26-input Process 구현과 Event 명세 검토가 끝난 뒤 수행한다.

\subsubsection{Recall과 Proceed의 최종 경계}

Transfer는 old finalized note를 소비하고 sender의 $cm_C$와 pending $rv$를 만든다.
Receiver용 finalized $cm$이나 $cm_{pending}$은 Transfer 시점에 만들지 않는다.

Proceed와 Recall은 모두 $rvnf=H(rv,o_{rv})$를 사용해 동일 voucher를 resolve한다.
Receiver는 unresolved voucher에 대해 Proceed를 실행해 fresh $cm_{recv}$를 만들 수 있다.
직전 sender는 $CurrentEpoch<\Map{rv}[rv]$인 동안 Recall을 실행해 fresh
$cm_{return}$을 만든다. Proceed는 Recall deadline이 지난 뒤에도 unresolved 상태라면
가능하며, protocol relation에는 deadline 상한 검사를 두지 않는다.

$o_{rv}$는 sender와 receiver가 공유하므로 두 당사자만 valid resolution witness를
구성할 수 있다. POC의 Event submitter는 permissionless이고 Sender/receiver
authorization은 각 Circuit의 owner relation으로 검증한다. RegisterUser에서 만든 EOA
mapping은 Event caller 제한에 사용하지 않는다. Public input 순서는 Protocol Specification의
$\mathbf{x}$ tuple에 적힌 왼쪽에서 오른쪽 순서를 따른다.

\subsubsection{Azeroth에서 재사용하는 경계}

Azeroth에서 재사용하는 것은 commitment, wallet note, Merkle membership, accepted root,
nullifier와 on-chain existence check의 skeleton이다. Azeroth의 단일 value $v$, encrypted
account, auditor encryption과 Public-Key Encryption을 그대로 가져오지 않는다. zkDPP는
quantity, StateVector와 8개 Event relation을 별도로 정의한다.

Transfer/Proceed에 필요한 note opening과 $o_{rv}$ 전달은 off-chain requirement로만
둔다. 구체적인 Public-Key Encryption scheme, ciphertext format과 encryption correctness
proof는 현재 POC와 논문의 main contribution에서 제외한다.

\paragraph{BLS12-381 Solidity 파일에서 참고하는 구조}

\texttt{BLS12\_381MixerBase.sol}, \texttt{Groth16BLS12\_381.sol}과
\texttt{Groth16BLS12\_381Mixer.sol}은 BLS12-381 Groth16 proof verification을
공통 Mixer에 연결하려던 과거 실험 경로다. 다음 설계 경계는 zkDPP에도 적용한다.

\begin{itemize}
  \item Protocol state와 business logic을 proof verifier에서 분리한다.
  \item 공통 contract는 root, nullifier, tree와 registry를 관리한다.
  \item Relation별 verifier는 proof와 ordered public input 검증만 담당한다.
  \item Curve-specific proof/point serialization과 precompile ABI를 verifier
        adapter 내부에 한정한다.
  \item Merkle tree는 abstract 2-to-1 hash 경계를 두고 Go, circuit과 Solidity가
        동일한 node compression을 사용한다.
\end{itemize}

\paragraph{재사용하지 않는 코드}

\begin{itemize}
  \item Azeroth BLS verifier는 Groth16용이며 zkDPP의 PLONK/KZG verifier가 아니다.
  \item 해당 verifier의 precompile address와 calldata layout은 최종 EIP-2537이 아닌
        과거 draft에 맞춰져 있으므로 복사하지 않는다.
  \item BLS Mixer 파일의 constructor와 \texttt{\_verifyZKProof} signature는 현재
        \texttt{ZklayBase}와 맞지 않으며 현 checkout의 compile/deploy target도 아니다.
  \item Azeroth \texttt{Poseidon2.sol}과 \texttt{MiMC7.sol}은 BN254 modulus와
        parameter를 사용하므로 BLS12-381 zkDPP hash로 재사용하지 않는다.
  \item Groth16 verifying-key storage layout, fixed input count, encrypted account,
        auditor encryption과 token/NFT branch를 가져오지 않는다.
\end{itemize}

zkDPP는 gnark가 export한 BLS12-381 PLONK Solidity verifier를 사용한다. Tree용
Poseidon2 Solidity compressor는 gnark-crypto v0.20.1 BLS12-381 default parameter에
맞춰 별도로 구현한다. 따라서 Azeroth에서 재사용하는 것은 architecture와 state-machine
skeleton이며 BLS verifier와 hash implementation source가 아니다.

\subsubsection{Poseidon2 결정 완료 범위}

Poseidon2에 관해 다음 항목은 모두 확정됐다.

\begin{itemize}
  \item Field: BLS12-381 scalar field
  \item Library: gnark v0.15.0과 gnark-crypto v0.20.1
  \item Default parameter: width 2, full rounds 6, partial rounds 50, S-box $x^5$
  \item Variable-length $H$와 $CM$: initial state 0의 Merkle--Damg\aa rd construction
  \item Protocol hash input ordering: 각 protocol 식의 왼쪽에서 오른쪽 flat ordering
  \item Poseidon2 해시 용도 구분값: 현재 POC에서는 입력에 추가하지 않음
  \item Merkle node hash: raw leaf에서 시작하는 direct 2-to-1 Poseidon2 compressor
  \item Cross-layer consistency: native Go와 gnark circuit은 제공 library를 재사용하고,
        Solidity 호환 구현은 동일 parameter와 reference vector로 검증
\end{itemize}

따라서 이전 절에 남은 ``Poseidon2 arity와 input ordering'' 질문은 이 후속 결정으로
대체된다. Reference vector 작성과 Solidity compressor 구현은 미확정 설계가 아니라
implementation task다.

DocumentHash의 Hash-to-Field primitive는 BLS12-381 \texttt{fr.Hash}로 확정됐다.
DocumentInfo encoding은 NFC, UTF-8와 unsigned 32-bit big-endian byte-length prefix로
확정됐다. \texttt{fr.Hash}의 DST는 빈 byte string이므로 별도 용도 구분값은 결정하지
않는다. 이는 off-chain DocumentHash encoding 문제이며 note commitment와 Merkle
Poseidon2 parameter를 다시 열지 않는다.

\subsubsection{Setup, Key와 Docker Foundry}

Entry, Transfer, Proceed, Recall, Merge, Split, Process와 Exit의 8개 Relation별로
PLONK proving/verifying key를 생성한다. POC에서는 각 Relation의 CCS를 gnark의
\texttt{unsafekzg.NewSRS(ccs)}에 전달해 canonical/Lagrange SRS pair를 만들고
PLONK Setup에 사용한다. 생성한 CCS, 두 SRS, proving/verifying key와 checksum manifest는
profile/Relation별 artifact directory에 저장하고, generated verifier는 Foundry source
tree의 profile/Relation directory에 저장한다. Manifest는 양쪽 file의 checksum을
binding한다. 이 SRS는
production SRS로 간주하지 않는다. Production ceremony와 integrity 검증은 Open Decision이다.

각 verifying key에서 export한 8개의 \texttt{PlonkVerifier}는 별도 contract로 배포한다.
Main zkDPP contract는 constructor에서 verifier address 8개를 받아 Relation ID와 binding하고,
공통 \texttt{Verify(bytes,uint256[])} ABI로 호출한다. 이를 ``verifier 배포 구조''라고 한다.
즉, proof 검증 code와 MT, nullifier, voucher state를 갱신하는 main protocol code를 서로
다른 contract로 나누고 연결하는 방식이다.

Azeroth 논문의 formal Scheme은 $pp=(ek,vk,G)$를 반환하고
\texttt{KeyGenUserClient(pp)}로 표기한다. 해당 KeyGen은 public-key encryption과
symmetric-key encryption의 key generation 및 Field sampling에 global parameter를
사용하지만 zk-SNARK proving key로 user ownership key를 만드는 것은 아니다.
Azeroth의 실제 Contract \texttt{registerUser}도 $pp$를 calldata로 받지 않고
off-chain에서 생성된 $addr$, $pkOwn$, $pkEnc$만 받는다. Formal algorithm의 $pp$ 전달과
Contract function argument는 서로 다른 abstraction level이다.

또한 local Azeroth checkout에는 Groth16 Setup/key generation tool이 없다. Deploy script는
\path{compiled/testParameter.json}의 사전 생성된 \texttt{vk}, \texttt{vk\_nft}를 읽어
main mixer constructor에 전달한다. 따라서 Azeroth는 zkDPP의
\texttt{unsafekzg.NewSRS} 선택을 정당화하는 구현 reference가 아니다.

Foundry는 host binary 설치를 protocol prerequisite로 두지 않는다. POC verifier test는
공식 Docker image \texttt{ghcr.io/foundry-rs/foundry:v1.7.1}을 사용하고
\texttt{evm\_version = "prague"}를 명시한다. 이 환경에서 EIP-2537 precompile을
사용하는 BLS12-381 PLONK verifier의 valid/invalid case를 확인했다.

\subsubsection{Entry와 사용자 등록의 경계}

Entry는 공급망 밖의 Product information을 새로운 finalized note로 발행하여 등록하는
$\varnothing\rightarrow cm_{new}$ Event다. Protocol은 DocumentHash별 최초 발행의
유일성을 강제하지 않는다. Entry 함수 안에서 address를 생성하지 않는다.
User는 먼저 KeyGen으로 $(usk,upk)$를 만들고 RegisterUser로 $addr$을 등록한 뒤, Entry가
이미 알고 있는 address를 commitment owner로 사용한다. 등록 단계의 address uniqueness,
EOA 재등록 거부와 $addr=H(pk_{own})$ Contract 검사는 Event ownership이나 gas overhead의
핵심 경로가 아니다. 이를 POC 필수조건으로 추가하면 아직 요구되지 않은 registry policy를
고정하므로 YAGNI 원칙에 따라 제거했다. POC mapping은 같은 EOA의 최신 값을 보관한다.

Azeroth 구현은 address 중복만 거부하고 $addr=H(pk_{own},pk_{enc})$를 Contract에서
검증하지 않는다. 또한 같은 EOA가 새로운 address로 다시 등록하면 기존 mapping을
덮어쓸 수 있으며, \texttt{LogUserRegister} event와 \texttt{getUserPublicKeys} getter를
모두 제공한다. zkDPP는 Azeroth의 registration policy를 그대로 복사하지 않으며,
stable EOA identity, key rotation과 registration event는 비차단 Open Decision으로 남긴다.

Entry issuer policy, direct
issuance 허용 여부와 origin data의 사실성은 아직 protocol이 보장하지 않는다.

\subsubsection{Smart Contract의 3단계 실행과 output 중복 처리}

각 Event 함수는 on-chain precondition 검사, proof verification, state update 순서로
실행한다. 기존 List 중복, accepted root, nullifier와 deadline처럼 proof 없이 확인할 수
있는 조건은 proof verification 전에 검사한다. Proof가 유효한 경우에만 List와 Map을
갱신하고 Tree에 append한다.

Split과 Process의 다중 output에는 별도 pairwise 비교를 추가하지 않는다. 모든 output을
proof 검증한 뒤 순서대로 기존 List에 없는지 검사하고 등록한 다음 Tree에 append한다.
뒤 output이 앞 output과 같으면 순차 check-and-insert가 실패한다. Solidity transaction의
atomicity에 의해 이 시점까지 수행한 nullifier 기록, List 등록, Tree node/root 변경과
event log는 모두 revert된다. 이 선택은 정상 transaction에 pairwise 비교 비용을 추가하지
않는 대신, 중복 output transaction이 proof verification gas를 소비한 뒤 실패하는
tradeoff를 수용한다.

\subsubsection{Merkle Path 조회와 Event 완결성 검토}

Client의 POC Merkle path 획득 방식을 Contract view 함수로 확정했다. Client는 Tree append
event에서 확인한 zero-based leaf index를 $MT$ 또는 $rvMT$의 path 조회 함수에 전달한다.
함수는 조회 시점의 current root와 leaf-to-root 순서의 sibling $d$개를 반환한다. Contract의
heap-indexed node mapping은 최신 node를 저장하므로 과거 root용 path를 재구성하지 않는다.
과거 root용 proof가 필요한 Client는 해당 시점의 path를 별도로 보관한다.

8개 Event의 Client, Relation과 Smart Contract를 다시 대조하면서 다음 구현 안전 조건을
명시적으로 고정했다.

\begin{itemize}
  \item Entry와 숫자 상태를 새로 계산하는 Transfer, Merge, Split, Process output의
        $q$와 StateVector coordinate는 unsigned 64-bit range를 검사한다.
  \item 기존 $MT$ 또는 $rvMT$ 객체를 소비하는 input은 생성 Relation에서 검증된 range를
        계승하므로 반복 검사하지 않는다. Proceed와 Recall의 Voucher 값 및 Exit input도
        같은 closed-state invariant를 적용한다.
  \item Merge의 $cm_1\neq cm_2$는 public input에 대한 사전 조건이므로 Circuit에서 제거하고
        Smart Contract가 proof verification 전에 검사한다.
  \item Merge의 두 input은 서로 다른 historical root를 사용할 수 있다. Public statement에
        $rt_1,rt_2$를 두고 $cm_i$의 Merkle path를 $rt_i$에 각각 binding하며, Contract는
        두 root가 모두 accepted $\List{rt}$에 있는지 검사한다.
  \item Split Client는 이미 준비된 $rt,path_{in}$을 외부 입력으로 받지 않고, wallet에
        저장된 $index_{in}$으로 current $rt$와 $path_{in}$을 조회한다.
  \item Exit Client도 같은 규칙을 사용해 $index_{in}$으로 current $rt$와
        $path_{in}$을 조회한다. Exit 관련 statement, witness, proof와 transaction 표기는
        축약형 $X$ 대신 $Exit$를 사용한다.
  \item Split은 zero-quantity output을 허용하지만 proportional floor/remainder와 aggregate
        conservation이 해당 output의 StateVector를 0으로 강제한다. $q_{in}=0$인 note는
        remainder bound를 만족할 수 없어 Split할 수 없다.
  \item 각 Smart Contract 함수는 transaction argument에서 Protocol Specification의
        ordered public statement $\mathbf{x}$를 직접 재구성한 뒤 verifier를 호출한다.
  \item Process의 active input commitment도 서로 달라야 한다. 그렇지 않으면 하나의
        finalized note를 여러 slot에 반복해 aggregate에 중복 반영할 수 있다.
  \item Split과 Process output은 별도 pairwise 비교를 추가하지 않는다. Proof 검증 후
        output별로 $\List{cm}$ 부재 확인, 등록과 Tree append를 순차 실행한다. 앞 output이
        먼저 등록되므로 같은 transaction 안의 중복도 기존 membership 검사로 거부된다.
  \item Process의 $m,n$은 별도 calldata 값이 아니라 active commitment 배열 길이에서
        Contract가 계산한다. Contract는 input root, commitment와 nullifier 배열 길이를
        비교하고 fixed verifier statement를 trailing zero로 padding한다.
  \item Process는 별도 $e^{in},e^{out}$ vector를 public input이나 Client witness로
        전달하지 않는다. Public $m,n$이 active prefix 길이를 나타내며 fixed circuit은
        $m,n$에 binding된 algebraic gating expression을 사용한다.
  \item Process의 각 input은 wallet에 저장된 $index_i$로 $(rt_i,path_i)$를 조회한다.
        Relation은 $cm_{in,i}$의 membership을 해당 $rt_i$에 binding하고 Contract는 active
        $rt_i$가 각각 accepted root인지 검사한다.
  \item Process Client와 transaction은 active $n\times\ell$ allocation matrix $A$만
        전달한다. Client는 proof 생성 전에, Contract는 verifier 호출 전에 각각 뒤의
        inactive row를 0으로 채워 동일한 fixed $N_{\max}\times\ell$ public input을 구성한다.
        이 선택은 verifier input 수를 줄이지 않지만 $n<N_{\max}$인 transaction calldata를 줄인다.
  \item Process Client는 output descriptor 배열
        $\{(\mathit{DocumentHash}_{out,j},q_{out,j})\}_{j=1}^{n}$을 입력받는다. $n$은 tuple
        cardinality로 이미 표현되므로 Client Scheme에서 배열 길이 계산과 arity assertion을
        반복하지 않는다. Process quantity와 DocumentHash transition을 강제하지 않으므로 이 두
        output 값은 $P$ 또는 $A$에서 도출하지 않는다.
  \item Split과 Process의 output 중복은 output pairwise constraint가 아니라
        proof 검증 후 실행되는 $\List{cm}$ 순차 등록으로 거부한다. 같은 transaction의
        뒤 output이 앞 output과 같으면 뒤 membership 검사에서 실패하고 transaction
        전체가 revert된다.
  \item Merkle node hash는 gnark-crypto v0.20.1 BLS12-381 Poseidon2
        \texttt{Compress} 의미론으로 고정했다. Permutation input은 $(left,right)$이고,
        결과의 오른쪽 lane에 원래 $right$를 더한 Field Element를 반환한다. Native,
        circuit과 Solidity는 동일 reference vector로 교차 검증한다.
\end{itemize}

이 항목들은 새로운 business policy가 아니라 기존 append-only set, 단일 소비와
fixed-slot 결정을 완전하게 구현하기 위한 안전 조건이다.

\subsubsection{Notation과 Scheme 작성 규칙}

Formal Scheme에서 $\mathbf{x}$와 $\mathbf{w}$는 ordered public statement와 private
witness다. 둘은 Azeroth식 중괄호로 표기하되 unordered set이 아니며, 왼쪽에서 오른쪽
순서를 보존하는 ordered tuple이다. Event 이름이 문맥에서 명확하므로 note, $\mathbf{x}$,
$\mathbf{w}$에 불필요한
Event suffix를 붙이지 않는다. $cm_{old}$와 $cm_{new}$는 상단 전이 요약에서 사용하고,
세부 Relation에서는 역할에 맞는 $cm_{in}$, $cm_C$, $cm_{recv}$ 등의 이름을 사용한다.

Private witness ordering은 $sk/pk$, DocumentHash와 object payload, address와 event 보조값,
Merkle path, opening/randomness 순서로 통일한다. Note와 voucher에 저장되는 field는 canonical
object 순서를 유지하고, 공유 field는 witness에 한 번만 둔다.

Voucher는 Sender용과 Receiver용 wallet record로 나누지 않는다. 다음 canonical
off-chain object 하나를 두 역할이 동일하게 저장한다.

\[
\mathit{voucher}=\{
rv,\mathit{DocumentHash},q_T,\mathbf{s}_T,
addr_{Sender},addr_{Receiver},\Delta_{epoch},o_{rv}
\}.
\]

각 participant가 자신의 address를 이미 알고 있어 일부 중복은 발생하지만, 역할별 schema와
복구 절차를 나누지 않는 일반성과 self-contained representation을 우선한다.

Proceed는 accepted $rvMT$ root에 대한 membership으로 voucher가 성공한 Transfer에서
생성됐음을 검증한다. 따라서 $\Map{rv}[rv]$의 별도 exists 검사를 하지 않는다. Recall도
$currentEpoch<\Map{rv}[rv]$ 비교 자체가 존재하지 않는 mapping의 기본값 0을 거부하므로
별도 exists 검사를 두지 않는다.

$rv\in\List{rv}$ 검사는 accepted $rvMT$ root와 Circuit membership proof가 이미 제공하는
membership soundness와 논리적으로 중복되지만 제거하지 않는다. Proceed와 Recall 모두
이 검사를 proof 검증 전 저비용 사전 filter로 유지한다. 별도의 $\Map{rv}$ exists 검사는
추가하지 않는다.

Process index는 input $i$, output $j$, State coordinate $k$, State length $\ell$로 고정한다.
이 표기 결정은 circuit array layout 자체를 확정하는 것이 아니라 논문, Scheme과 구현
문서가 같은 차원을 가리키도록 하기 위한 것이다.

\subsubsection{최신 Protocol 구현을 위한 새 POC 분리}

\texttt{Protocol Spec.tex}을 최신 canonical Protocol 문서로 확정한다.
기존 \path{poc/}와 저장소 root의 \path{contracts/}는 Entry와 고정 3-to-3 Process의
과거 성능 결과를 재현하는 historical baseline으로 보존하며 최신 Protocol에 맞게
수정하지 않는다. 최신 8-Event 구현은 독립된 \path{poc-v2/}에서 진행한다.

기존 POC의 Event circuit, main Contract, fixture, generated verifier와 benchmark 결과는
이전 public statement에 binding되어 있으므로 재사용하지 않는다. 반면 dependency version,
검증된 Poseidon2 hash/compressor, Depth-32 Merkle Tree, Solidity Poseidon2 generator,
artifact writer와 Docker/Foundry 환경은 Protocol 의미를 포함하지 않는 공통 engineering
component이므로 reference vector와 단위 test를 다시 통과하는 조건으로 복사해 사용한다.
새 POC는 기존 module을 import하지 않으며 key, verifier, fixture와 artifact directory도
공유하지 않는다.

\subsubsection{Canonical scenario와 benchmark 구현 결과}

Canonical scenario의 Entry State는 모두 $10^9$의 배수인 whole-unit 값으로 고정했다.
반면 Split, partial Transfer와 Process에는 nonzero deterministic remainder를 유지하여
first-output residual relation을 실제 proof에서 검증한다. 사람이 결정한 원본 값은
\path{poc-v2/testdata/scenario/canonical.json}에 한 번만 기록하며, hash, commitment,
nullifier, voucher, Merkle root와 remainder는 derived artifact로 자동 생성한다.

최신 POC는 8개 Event의 Go proof, generated Solidity verifier와 Main Contract state update를
canonical scenario에서 모두 통과했다. Benchmark는 SHA와 depth scaling을 제외하고
Poseidon2, depth 32의 34개 distinct compile-time profile로 확정했다. Profile별 verifier
gas와 canonical Event gas는 측정 경계가 다르므로 하나의 수치로 합치지 않고 별도 artifact로
기록한다.

\subsubsection{Policy/Recipe 기반 Process의 검토와 POC 제외}

현재 generic Process는 산업 공정 전체를 일반화하지 않는다. Input State를 aggregate한
뒤 공개된 process delta $P$와 StateAllocationMatrix $A$를 적용하여 output State가
선언된 식과 일치하는지만 증명한다. Entry issuer가 초기 State를 정직하게 등록하고
Process operator가 실제 공정을 반영한 $P,A$를 공개한다는 가정이 남는다. 따라서 public
$P,A$는 악의적인 선언을 암호학적으로 금지하지 않고, 선언을 숨기지 못하게 하여 외부
검토와 사후 책임 추적을 가능하게 한다.

주기적인 사후 감사에 의존하지 않으려면 Process가 허용되는 input/output Material class,
Quantity transition rule, State delta policy와 State allocation policy를 Recipe로 직접
검증해야 한다. 이를 구현하는 큰 선택지는 Recipe별 circuit을 두는 방식과, 하나의 공통
circuit이 승인된 Recipe를 private witness로 검증하는 방식이다. Recipe별 circuit은
verifier 선택으로 공정과 산업을 노출할 수 있다. Hidden Recipe commitment와 membership
proof 방식은 Recipe 승인 주체, registry governance, Merkle data availability와 path
획득 문제를 새롭게 만든다. Permissionless registry는 Recipe의 산업적 정당성을 보장하지
못하고, 중앙 승인자는 기존 supply-chain audit를 다른 Trusted Party로 옮기는 문제가 있다.

따라서 이번 POC에는 \texttt{MaterialClassId}, \texttt{RecipeId}, Recipe commitment,
Recipe Merkle Tree/root, Recipe membership proof와 Recipe별 verifier를 추가하지 않는다.
기존 public $P,A$ State Allocation Process를 먼저 검토하고 구현한다. Policy/Recipe
Process는 trust model, Recipe privacy와 승인 구조를 함께 정의할 수 있을 때 후속
연구로 다시 검토한다.

\end{document}
